\chapimg{fig/op03.png}
\IfFileExists{extra/lead03.tex}{\chapterlead{\subhd{After this chapter you can} classify irrigation triggers and sensors; compute the depletion that a suction trigger means and the suction that a depletion means; derive and evaluate the reach of a trigger; apply the dose law for the fate of an irrigation millimetre; define and use the memory fate of a flux; and choose a trigger for a soil in a design.}
}{}
\chapter{Triggers: Reach, Fate and Hydraulic Memory}\label{chap:03}
\markboth{Triggers: Reach, Fate and Hydraulic Memory}{Triggers: Reach, Fate and Hydraulic Memory}
\IfFileExists{extra/bg03.tex}{\section{Introduction: when to irrigate}
Every irrigation system answers two questions: \emph{when} to apply water and \emph{how much}. The second is a matter of the root zone and the system (Chapters~2 and~7); the first is a matter of a \kw{trigger}, a rule that turns the state of the soil, the crop or the weather into a decision. This chapter studies the simplest trigger, a threshold on the suction measured by a sensor in the root zone, and asks three questions of it: how far can a soil dry before the trigger fires (the \emph{reach}), where does the water go after it is applied (the \emph{fate}), and how long does the soil remember a change of management (the \emph{memory}). The answers are laws derived in the studies on which this book is based from an ensemble of 10\,190 Daisy simulations over ten soils.

\section{Types of irrigation triggers and sensors}
Irrigation scheduling methods are grouped into four classes. \kw{Calendar} (fixed interval or fixed dose) methods ignore the state of the crop and are the oldest; they waste water when it rains and starve the crop in a hot spell. \kw{Water balance} (or budget) methods keep a running account of the depletion $D_r$ of the root zone from the weather ($ET_c$) and the water applied, and irrigate when $D_r$ reaches the readily available water; they need no sensor but drift without correction. \kw{Soil-based} methods read a sensor: a \emph{tensiometer} (a porous cup on a water-filled tube with a vacuum gauge, reading 0--80 kPa), a \emph{gypsum block} or \emph{granular matrix sensor} (electrical resistance, 0--200 kPa), or a \emph{dielectric probe} (water content). \kw{Plant-based} methods read leaf water potential, canopy temperature or stem diameter.

Soil-based triggers are of two kinds: a \emph{suction trigger} $h_t$ (irrigate when the sensor reads more than $h_t$ cm) and a \emph{depletion trigger} $f_t$ (irrigate when a fraction $f_t$ of the available water has gone). The two are linked by the retention curve of the soil, $f=f(h)$ (Equation~\ref{eq:fdep}, Chapter~9). A suction of 300 cm is a light depletion in a sand and a deep one in a silt loam; that is the reason the same trigger behaves differently in different soils.

\exbox{A suction trigger as a depletion}{For the sand S03 ($\alpha=0.145$, $n=2.68$) and the silt loam S07 ($\alpha=0.020$, $n=1.41$) find the depletion $f$ that a suction trigger of 300 cm means (field capacity at 100 cm, wilting point at 15\,849 cm).}\sol \givenline{$\alpha=0.145$ ; $n=2.68$ ; $\alpha=0.020$ ; $n=1.41$}\ans Sand: $f(300)=[S_e(100)-S_e(300)]/[S_e(100)-S_e(15849)]=(0.0112-0.0018)/(0.0112-0.0000)=0.842$; Silt loam: $f(300)=[S_e(100)-S_e(300)]/[S_e(100)-S_e(15849)]=(0.6859-0.4691)/(0.6859-0.0943)=0.366$ \ans The same 300 cm trigger is a near-complete depletion in the sand and a modest one in the silt loam. \finans

\section{The reach of a trigger}
A trigger can fire only if the soil can reach the suction at which it is set. A crop dries the soil until the conductivity around its roots falls so low that the soil can no longer supply the demand; beyond that suction the roots meet a wall. Writing that $K_r(h_c)=\kappa$ at that point, and using the dry-end form $K_r\simeq m^2(\alpha h)^{-E}$ with $E=2n+\ell(n-1)$ (Chapter~4, Equations on the reach), the reach is

\begin{equation}h_c=\frac1\alpha\Bigl(\frac{m^2}{\kappa}\Bigr)^{1/E},\qquad m=1-\frac1n,\ \kappa=3.5\times10^{-10},\ \ell=0.5.\label{eq:reach3}\end{equation}
\par\noindent where $h_c$ is the reach (cm), $\alpha$ (cm$^{-1}$) and $n$ the van Genuchten parameters, $m=1-1/n$, $\kappa$ the demand constant, $\ell$ the connectivity parameter and $E=2n+\ell(n-1)$ the exponent.



A trigger $h_t<h_c$ can fire and irrigates; a trigger $h_t>h_c$ never fires and the field behaves as rainfed. For a design, this gives an upper limit for the trigger of each soil.

\begin{table}[!htbp]\centering\small\caption{Reach of a suction trigger for the ten soils ($\kappa=3.5\times10^{-10}$, $\ell=0.5$) computed from the rounded parameters.}
\begin{tabular}{llrrrr}\toprule
\hd{Soil} & \hd{Class} & \hd{$n$} & \hd{$E$} & \hd{$h_c$ (cm)} & \hd{300 cm fires?}\\\midrule
S01 & Coarse loam & 1.35 & 2.88 & 7\,814 & yes\\
S02 & Coarse sand & 1.60 & 3.50 & 5\,359 & yes\\
S03 & Sand & 2.68 & 6.20 & 199 & no\\
S04 & Loamy sand & 2.28 & 5.20 & 425 & yes\\
S05 & Sandy loam & 1.89 & 4.22 & 1\,615 & yes\\
S06 & Loam & 1.56 & 3.40 & 9\,186 & yes\\
S07 & Silt loam & 1.41 & 3.02 & 29\,528 & yes\\
S08 & Sandy clay loam & 1.48 & 3.20 & 7\,559 & yes\\
S09 & Clay loam & 1.31 & 2.78 & 47\,608 & yes\\
S10 & Silty clay loam & 1.23 & 2.58 & 127\,830 & yes\\\bottomrule\end{tabular}\end{table}

\exbox{Reach of a loam}{Compute the reach of the loam S06 ($\alpha=0.036$, $n=1.56$) step by step, and say whether a trigger of 1000 cm would fire.}\sol \givenline{$\alpha=0.036$ ; $n=1.56$}\ans $m=1-1/1.56=0.3590$, $m^2=0.1289$; $E=2n+\ell(n-1)=3.12+0.5\times0.56=3.40$. $m^2/\kappa=0.1289/3.5\times10^{-10}=3.682e+08$; its power $1/E=0.2941$ is $330.70$, so $h_c=330.70/0.036=9186$ cm \ans A trigger of 1000 cm is below the reach and fires. \finans

\section{The fate of an irrigation millimetre: the dose law}
An irrigation millimetre does not all go to the crop. Written against a rainfed twin of the same run, the water applied $I$ is shared among transpiration, soil evaporation, other evaporation, drainage and storage; the shares $\eta_X$ (the \kw{marginal fate} in flux $X$) sum to one. The drained share depends mostly on the dose $d$ per event, because a larger dose wets deeper and exceeds the capacity of the root zone more often:

\begin{equation}\phi_D(d)=\eta_D\bigl(1-e^{-d/d_D}\bigr),\label{eq:dose3}\end{equation}


where $\eta_D$ is the saturation level of the drained share and $d_D$ the dose scale. The law is a fit to the ensemble, not a derivation; its use is to choose a dose that keeps drainage low.

\exbox{Drained share against dose}{A soil with $\eta_D=0.45$ and $d_D=9$ mm gives the dose law. Find the drained share at doses of 5, 10, 15, 20 and 30 mm per event, and the dose at which half of the saturation level is reached.}\sol \givenline{$\eta_D=0.45$ ; $d_D=9$ mm gives}\ans $d=5$: $\phi_D=0.45(1-e^{-5/9})=0.192$; $d=10$: $\phi_D=0.45(1-e^{-10/9})=0.302$; $d=15$: $\phi_D=0.45(1-e^{-15/9})=0.365$; $d=20$: $\phi_D=0.45(1-e^{-20/9})=0.401$; $d=30$: $\phi_D=0.45(1-e^{-30/9})=0.434$. Half of saturation: $1-e^{-d/9}=0.5$, $d=9\ln2=6.2$ mm \ans Doses below about 6 mm lose less than a third of the saturation level to drainage, so frequent small doses are efficient in this soil; a 30 mm dose loses almost the full 0.45. \finans

\section{Hydraulic memory}
When irrigation is changed (a mulch is added, a trigger lowered), the soil does not respond at once, because the water that was held before the change must drain or be used. This persistence is the \kw{hydraulic memory}. If a saving of $s$ millimetres is made in a flux in year~1, the \kw{memory fate} $\phi_x$ of a flux $x$ is the share of the saving that appears in $x$ in the following years; it behaves as a pooled slope of the flux change on the storage change. Large in soils that hold water (fine soils) and small in free-draining soils, it limits how much of a saving can be ``banked'' in the root zone. Chapter~7 of the application tabulates it for ten soils (Table at the end of this chapter).

}{}
\banner{The reach of a suction trigger}\label{ch:reach}
\section{The question}
Sensor-scheduled irrigation applies water when the suction at a sensor passes a trigger $h_t$. Triggers of 300 to 600 cm come from tables, the same for every soil. But a soil may never reach such a suction. If the crop cannot dry the soil beyond some suction $\hc$, then a trigger $h_t>\hc$ is never reached, the crop is not irrigated, and the field is in effect rainfed. A trigger $h_t<\hc$ can fire. The quantity $\hc$ is the \emph{reach} of the trigger.

\begin{kfig}{Sensor suction through a season: the days beyond the trigger $h_T$ are counted; rainfed (dashed) crosses the trigger earlier and stays beyond it for longer than irrigated.}
\draw[axis] (0,0)--(7.6,0);\draw[axis] (0,0)--(0,4.2);
\node[lab,below] at (3.8,-0.45) {Day of the season};\node[lab,rotate=90] at (-0.8,2.1) {Suction $h$ (cm)};
\draw[gd] (0,2.6)--(7.4,2.6);\node[lab,right] at (7.4,2.6) {$h_T$};
\draw[crv,smooth] plot coordinates {(0,0.4)(1,0.6)(2,1.2)(3,2.0)(3.5,2.75)(3.6,1.0)(4.4,1.4)(5.2,2.0)(5.9,2.8)(6.0,1.0)(6.8,1.5)(7.4,2.1)};
\draw[ln,dashed,smooth] plot coordinates {(0,0.4)(1,0.9)(2,1.8)(3,2.7)(4,3.3)(5,3.7)(6,3.9)(7.4,4.0)};
\fill[khcyan!25] (3.0,2.6) rectangle (3.5,2.75);\node[lab] at (1.4,3.6) {irrigated: ---};\node[lab] at (1.5,3.2) {rainfed: - - -};
\end{kfig}

\section{The dry-end asymptote}
Far from saturation, $\alpha h\gg1$. Then $\Se^{1/m}=[1+(\alpha h)^n]^{-1}\approx(\alpha h)^{-n}$ and, since $m n=n-1$, $\Se\approx(\alpha h)^{-(n-1)}$. In Equation~(\ref{eq:vgm}), $1-(1-y)^m\approx m y$ for small $y=\Se^{1/m}$, so the bracket is $\approx m(\alpha h)^{-n}$ and
\begin{equation}\label{eq:asym}\boxed{\boxed{\;\Kr(h)\simeq m^{2}(\alpha h)^{-[\,2n+\ell(n-1)\,]},\qquad\alpha h\gg1.\;}}\end{equation}
\par\noindent where $K_r$ is the relative conductivity, $h$ the suction (cm), $\alpha$, $n$ and $m$ the van Genuchten parameters, and $\ell$ the connectivity parameter; the power-law form holds when $\alpha h\gg1$, that is in dry soil.

The conductivity is a power law of suction, and the exponent $E=2n+\ell(n-1)$ grows with $n$.

\exbox{How good is the asymptote?}{For the silty clay loam S10 ($\alpha=0.010$, $n=1.23$, $m=0.187$, $E=2.575$): at $h=1000$ cm the exact value is $\Kr=8.64\times10^{-5}$ and Equation~(\ref{eq:asym}) gives $9.30\times10^{-5}$, an error of 8 per cent. At $h=10{,}000$ cm the exact value is $2.46\times10^{-7}$ and the asymptote $2.48\times10^{-7}$, an error under 1 per cent. At 100 cm, where $\alpha h=1$, the asymptote $3.5\times10^{-2}$ is wrong by a factor of 2.5, because the assumption $\alpha h\gg1$ fails. The asymptote is therefore reliable at the suctions where triggers sit.}



\section{One constant for all soils}
Suppose the crop stops drying the soil when the relative conductivity has fallen to a value $\kappa$. The justification is Chapter 4: the uptake collapses once $K$ is small enough. Setting $\Kr(\hc)=\kappa$ in Equation~(\ref{eq:asym}),
\begin{equation}\label{eq:reach}\boxed{\boxed{\;\hc=\frac1\alpha\left(\frac{m^{2}}{\kappa}\right)^{1/[2n+\ell(n-1)]}\;}}\end{equation}
\par\noindent where $h_c$ is the reach, the suction (cm) beyond which the soil can no longer supply the demand of the crop, $\alpha$, $m$, $n$ and $\ell$ the van Genuchten parameters, $\kappa$ a dimensionless constant that lumps the demand of the crop, the climate and the geometry of the roots (here $3.5\times10^{-10}$ in the application), and $E=2n+\ell(n-1)$ the exponent of the power law.

The constant $\kappa$ describes the crop, the climate and the sensor. It does not depend on the soil, which enters only through $\alpha$, $n$ and $\ell$. With $\kappa=3.5\times10^{-10}$ and $\ell=0.5$:

\begin{table}[!htbp]\centering\small\caption{Reach of a trigger from Equation~(\ref{eq:reach}), $\kappa=3.5\times10^{-10}$, $\ell=0.5$, computed from the rounded parameters of Table~\ref{tab:soils}; the study's values differ by under 1 per cent. The last column: does a 300 cm trigger fire?}\label{tab:reach}
\begin{tabular}{llrrrrl}\toprule
\hd{ID} & \hd{Texture} & \hd{$n$} & \hd{$E$} & \hd{$\alpha\hc$} & \hd{$\hc$ (cm)} & \hd{300 cm}\\\midrule
S03&Sand&2.680&6.200&28.8&199&never\\S04&Loamy sand&2.280&5.200&52.7&425&never\\S05&Sandy loam&1.890&4.225&121&1,615&fires\\S02&Coarse sand&1.600&3.500&287&5,359&fires\\S08&Sandy clay loam&1.480&3.200&446&7,559&fires\\S01&Coarse loam&1.353&2.882&751&7,814&fires\\S06&Loam&1.560&3.400&331&9,186&fires\\S07&Silt loam&1.410&3.025&591&29,528&fires\\S09&Clay loam&1.310&2.775&905&47,608&fires\\S10&Silty clay loam&1.230&2.575&1,278&127,830&fires\\\bottomrule\end{tabular}\end{table}
(The study reports values above 10,000 cm as bounded from below, because the trigger experiment stops at 10,000 cm.)

\exbox{The reach of a soil by hand}{A soil has $\alpha=0.05$ cm$^{-1}$ and $n=2.0$. Then $m=0.5$, $E=2(2.0)+0.5(1.0)=4.5$, $m^2/\kappa=0.25/(3.5\times10^{-10})=7.14\times10^{8}$, $\ln(7.14\times10^8)=20.39$, and $20.39/4.5=4.531$, so $(m^2/\kappa)^{1/E}=e^{4.531}=92.9$. The reach is $\hc=92.9/0.05=1{,}856$ cm. A trigger of 300 cm fires in this soil. In practice one keeps the trigger below about half the reach, here under 930 cm, because irrigation falls to zero as $h_t/\hc$ approaches 1 (figure at the end of this chapter).}



\subsection{Properties of the reach}
\begin{enumerate}[leftmargin=*,itemsep=2pt]
\item $\alpha$ only rescales the reach: $\alpha\hc$ depends on $n$, $\ell$ and $\kappa$ alone.
\item \emph{Sensitivity to $n$.} Write $\ln(\alpha\hc)=\ln(m^2/\kappa)/E$. With $dm/dn=1/n^2$,
\begin{equation}\label{eq:sens}\boxed{\frac{\partial\ln(\alpha\hc)}{\partial n}=\frac{2}{m\,n^2E}-\frac{(2+\ell)\ln(m^2/\kappa)}{E^2}.}\end{equation}
\par\noindent where $\partial\ln(\alpha h_c)/\partial n$ is the logarithmic sensitivity of the reach to the shape parameter $n$; $m=1-1/n$, $E=2n+\ell(n-1)$, $\ell$ and $\kappa$ are as in Eq.~(\ref{eq:reach}).

The second term dominates, since $\ln(m^2/\kappa)$ is between 18 and 21. The reach therefore shortens steeply as $n$ grows: $\alpha\hc\approx1{,}280$ at $n=1.23$ and $\approx29$ at $n=2.68$.
\item A larger $\ell$ shortens the reach, because the conductivity collapses faster.
\item For $n\to1$ the exponent tends to 2 and $m\to0$, so the reach tends to $\alpha^{-1}(m^2/\kappa)^{1/2}$, which vanishes slowly. Very broad pore-size distributions remain reachable far into the dry range, as observed for the silty clay loam.
\item \emph{Robustness.} Because $\hc\propto\kappa^{-1/E}$, a tenfold change of $\kappa$ changes the reach by $10^{1/E}$: 1.45 in the sand ($E=6.2$) and 2.44 in the silty clay loam ($E=2.575$).
\end{enumerate}

\begin{keybox}A suction trigger has a reach $\hc=\alpha^{-1}(m^2/\kappa)^{1/[2n+\ell(n-1)]}$. A trigger beyond it can never fire. One constant, $\kappa\approx3.5\times10^{-10}$, serves all ten soils; left-out soils are predicted within a factor of 2.3.\end{keybox}

\section{Why a trigger has a reach, in words}
A crop dries the soil at the sensor only as fast as the soil delivers water to the roots. That delivery scales with the conductivity, which falls by many orders of magnitude as the soil dries. In a soil with a narrow pore-size distribution (large $n$, sands) the conductivity collapses within a narrow band of suction just past the air entry, so the crop stops drying the soil there, and whatever trigger is set beyond that band is never read. In a soil with a broad distribution (small $n$, silty clay loams) the conductivity fades slowly over many decades, and the crop can still dry the soil at very high suctions.

\section{Why a deeper sensor does not help}
Below the sensor depth of the experiments the root density is smaller, so the soil dries even less, and in sands the deeper soil is also wetted by drainage from above. A deeper sensor therefore reads lower suctions, not higher ones.

\section{The reach scale}
Define the reach-scaled trigger $x=h_t/\hc$. Plotting irrigation against $x$ instead of against $h_t$ collapses the curves of ten soils: the mean within-bin variance of the normalised irrigation, in ten quantile bins, falls from 54 per cent of the total variance (on $h_t$ alone) to 32 per cent (on $x$). It does not fall to zero, because the reach sets where irrigation ends, but how irrigation fades \emph{before} the reach depends on how often the soil reaches a given suction in a season, which depends on the water held between field capacity and that suction: a retention property, not a conductivity property.

\label{fig:reachscale}

\exbox{}{Compute $\hc$ for $\alpha=0.03$ cm$^{-1}$, $n=1.7$ with $\kappa=3.5\times10^{-10}$ and $\ell=0.5$.}\sol  $m=0.4118$, $E=2(1.7)+0.5(0.7)=3.75$, $\ln(m^2/\kappa)=20.00$, so $(m^2/\kappa)^{1/E}=e^{5.333}=207$ and $\hc=207/0.03\approx6{,}900$ cm.

\exbox{}{Derive Equation~(\ref{eq:asym}) from Equation~(\ref{eq:vgm}), step by step, stating each approximation.}\sol  See the derivation in Section 5.2.

\exbox{}{Derive Equation~(\ref{eq:sens}). Evaluate it at $n=1.5$ with $\kappa=3.5\times10^{-10}$ and $\ell=0.5$ and say whether the reach rises or falls with $n$.}\sol  At $n=1.5$: $m=0.3333$, $E=3.25$, $\ln(m^2/\kappa)=19.57$; first term $2/(0.3333\cdot2.25\cdot3.25)=0.820$; second $2.5\cdot19.57/10.56=4.633$; derivative $-3.81$: the reach falls steeply with $n$.

\exbox{}{For the soil of the worked example, by what factor does the reach change if $\kappa$ is $3.5\times10^{-9}$ instead of $3.5\times10^{-10}$?}\sol  A factor of $10^{1/E}=10^{1/4.5}=1.67$ shorter for $\kappa$ ten times larger ($1{,}856\to1{,}113$ cm).

\exbox{}{A farmer sets a 300 cm trigger in a loamy sand. Use Table~\ref{tab:reach} to explain why the field is in effect rainfed, and propose a trigger.}\sol  The reach of the loamy sand is 425 cm and 300 cm is $0.7\hc$, so irrigation is nearly nil; keep the trigger below $\hc/2\approx210$ cm.
\banner{The fate of an irrigation millimetre}\label{ch:fate}
\section{Marginal, not average}
Of every millimetre of water pumped onto a field, how much does the crop transpire, and where does the rest go? The answer must hold for many soils, doses, seasons and methods. Take an irrigated run and its \emph{rainfed twin}, which has the same soil, weather and crop. The long-term balances are
\begin{equation}\label{eq:twin}\boxed{P+I=T+E_s+E_o+D+\Delta S+Q,\qquad P=T_0+E_{s,0}+E_{o,0}+D_0+\Delta S_0+Q_0 .}\end{equation}
\par\noindent where $P$ is the precipitation, $I$ the irrigation, $T$ the transpiration, $E_s$ the soil evaporation, $E_o$ the other evaporation (canopy and mulch), $D$ the drainage, $\Delta S$ the change of stored water and $Q$ the runoff (all in mm over the season); the subscript 0 marks the rainfed twin that receives no irrigation.

Subtracting and dividing by $I$,
\begin{equation}\label{eq:fate}
\boxed{\;1=\underbrace{\frac{T-T_0}{I}}_{\eta_T}+\underbrace{\frac{E_s-E_{s,0}}{I}}_{\eta_{E_s}}+\underbrace{\frac{E_o-E_{o,0}}{I}}_{\eta_{E_o}}+\underbrace{\frac{D-D_0}{I}}_{\eta_D}+\underbrace{\frac{\Delta S-\Delta S_0+Q-Q_0}{I}}_{\eta_S}\;}
\end{equation}
\par\noindent where $\eta_X=(X-X_0)/I$ is the share of the irrigation $I$ (mm) that appears as the change of the term $X$ ($X=T,E_s,E_o,D$, or storage plus runoff) relative to the rainfed twin (subscript 0); the five shares add to one because the balance of Eq.~(\ref{eq:twin}) closes.

Each $\eta_X=(X-X_0)/I$ is the share of the irrigation that appears as a change in flux $X$. The shares are \emph{marginal}: they describe the response to the water that was added, not the share of the total in the irrigated run, so a share can be negative (overhead irrigation can lower transpiration) or larger than the average loss fraction. Pooling a group of runs uses the ratio of sums, $\eta_X=\sum(X-X_0)/\sum I$, which weights each run by its irrigation.

\exbox{One pair of runs}{An irrigated run has $I=200$ mm, and compared with its rainfed twin: $T-T_0=+76$, $E_s-E_{s,0}=+50$, $E_o-E_{o,0}=+38$, $D-D_0=+32$ and $\Delta S-\Delta S_0+Q-Q_0=+4$ mm. These add to $200$ mm. The shares are $\eta_T=0.38$, $\eta_{E_s}=0.25$, $\eta_{E_o}=0.19$, $\eta_D=0.16$ and $\eta_S=0.02$ --- the average fate found over the ensemble. Less than half of the water is transpired.}



\fig{p4_fate}{0.9}{The fate of an irrigation millimetre by soil and method, from the Daisy ensemble.}

\section{What the ensemble shows}
Over eight soils, ten doses, ten seasons and three methods, each irrigation millimetre became 0.38~mm of transpiration, 0.25~mm of soil evaporation, 0.19~mm of other evaporation, 0.16~mm of drainage below 1~m and 0.02--0.03~mm of storage and runoff. The transpired share is largest in sandy loam (0.64) and smallest in coarse sand (0.23), where 0.32 drains. In the five fine-textured soils the shares are close: 0.34 to 0.46 transpired and 0.15 drained.

\subsection{Overhead irrigation adds no transpiration}
In the methods experiments the extra transpiration of overhead irrigation is nil or negative ($-0.08$ in E4 and $-0.01$ in E15): the water that the canopy intercepts evaporates (0.79 and 0.56 of each millimetre) and lowers the transpiration that it replaces. Subsurface application transpires most (0.51 in E4, 0.40 in E15) but also drains most. Why: transpiration is limited by the energy available to the canopy and by the cover; water applied to a wet surface or to a canopy that does not cover the ground evaporates using the same energy, and water applied beyond the storage of the root zone drains.

\begin{kfig}{The drained share of the water applied rises with the dose per event $d$ and saturates at $\eta_D$: $\phi_D=\eta_D\,(1-e^{-d/d_D})$.}
\draw[axis] (0,0)--(7.4,0);\draw[axis] (0,0)--(0,4.2);
\node[lab,below] at (3.7,-0.55) {Dose per event $d$ (mm)};\node[lab,rotate=90] at (-0.9,2.1) {Drained share $\phi_D$};
\foreach \x/\t in {0/0,1.5/10,3/20,4.5/30,6/40}{\draw[ln] (\x,0)--(\x,-0.1);\node[lab,below] at (\x,-0.1) {\t};}
\draw[gd] (0,3.4)--(7.2,3.4);\node[lab,above right] at (0.1,3.4) {$\eta_D$};
\draw[crv,domain=0:7,samples=50,smooth] plot (\x,{3.4*(1-exp(-\x/1.9))});
\draw[gd] (1.9,0)--(1.9,2.149);\draw[gd] (0,2.149)--(1.9,2.149);\node[lab,below] at (1.9,-0.5) {$d_D$};
\end{kfig}

\section{The dose law}
The drained share rises with the dose $d$ applied per event and levels off. A simple form with the right limits is
\begin{equation}\label{eq:dose}\boxed{\boxed{\;\eta_D(d)=\eta_\infty\bigl(1-e^{-d/d_D}\bigr)\;}}\end{equation}
\par\noindent where $\eta_D$ is the share of the irrigation that drains (dimensionless), $d$ the dose applied per event (mm), $\eta_\infty$ the limit of the share for very large doses and $d_D$ the characteristic dose (mm) at which 63 per cent of the limit is reached.

\textbf{Why this form.} Let a dose be applied when the root zone has a deficit. Part of the dose is used by the crop and the surface before it can move below the roots; the rest can drain. If the water usable before redistribution grows with the dose and saturates, and the drained share of the remainder is a soil property $\eta_\infty$, then Equation~(\ref{eq:dose}) is the simplest form with the right limits: $\eta_D\to\eta_\infty d/d_D$ for $d\ll d_D$ (a vanishing dose cannot drain) and $\eta_D\to\eta_\infty$ for $d\gg d_D$. The dose scale $d_D$ is the dose at which $1-e^{-1}=63$ per cent of the plateau is reached. \emph{The law is phenomenological}: it was chosen for its limits and tested on data (coefficient of determination 0.92 to 0.96 on seven soils), not derived from the Richards equation.

\subhd{Findings.} The dose scale is $d_D=4.1\pm0.6$ mm, the same in every soil and close to one day of crop water use. The plateau $\eta_\infty$ is soil specific, between 0.10 and 0.45, and highest in coarse sand.

\subhd{Why the scale is the same in every soil.} The water that a dose can lose to drainage is the part that the crop and the surface do not use before it redistributes below the roots. The crop uses about one day of evapotranspiration in that time whatever the soil, so the dose scale is set by the atmosphere and the crop, not by the soil. \textbf{Why the plateau depends on the soil:} once a dose is large, the share that drains is the share of the excess that the soil passes on before the crop uses it, a race between drainage and uptake. Coarse sand passes water fast, fine soils hold it.

\exbox{Choosing a dose}{Take a soil with $\eta_\infty=0.30$ and $d_D=4.1$ mm. Equation~(\ref{eq:dose}) gives $\eta_D=0.065$ at $d=1$ mm, $0.19$ at 4.1 mm, $0.257$ at 8 mm, $0.284$ at 12 mm and $0.298$ at 20 mm. Increasing the dose from 4.1 to 12 mm raises the drained share from 0.19 to 0.28, that is, 9 extra drained millimetres per 100 mm applied. But a dose of 1 mm drains almost nothing and loses 60 to 72 per cent to evaporation, so very small doses are no solution either: the transpired share rises with the dose up to 15 to 20 mm in most soils.}



\fig{p4_dose}{0.9}{Drained, transpired and evaporated share against the dose per event, with the fit of Equation~(\ref{eq:dose}).}

\begin{daisybox}{computing the fate and the dose law}
For each soil and each run of E10 (dose 1 to 30 mm applied in 12 h, ten irrigation windows), the annual drainage at 1 m is read from the Daisy log and subtracted from the drainage of the rainfed twin in E0. Pooling gives $\eta_D$ per soil and dose. Equation~(\ref{eq:dose}) is then fitted per soil by nonlinear least squares for $\eta_\infty$ and $d_D$. Methods come from E4 (surface, overhead, subsurface on each soil) and from E15 (method crossed with rate, dose and rooting depth in coarse loam). All shares sum to one by construction, so the fifth is obtained as a check, not fitted.
\end{daisybox}

\exbox{}{A run has $I=150$ mm and differences from its twin of $T$: $+45$, $E_s$: $+40$, $E_o$: $+30$, $D$: $+30$. Find $\eta_S$ and all five shares.}\sol  $\eta_S=150-(45+40+30+30)=5$ mm, so shares $0.30,\ 0.267,\ 0.20,\ 0.20,\ 0.033$.

\exbox{}{Show that a negative $\eta_T$ is possible when the added water is intercepted and evaporated.}\sol  If the added water is intercepted and evaporates, canopy evaporation of the irrigation can exceed what the crop would otherwise take up, so the change of $T$ can be below zero.

\exbox{}{For $\eta_\infty=0.45$ and $d_D=4.1$ mm compute $\eta_D$ at 2, 5, 10 and 25 mm. At what dose does $\eta_D$ reach 90 per cent of the plateau?}\sol  $\eta_D=0.45(1-e^{-d/4.1})$: $0.147$, $0.250$, $0.411$, $0.447$ at 2, 5, 10, 25 mm. 90 per cent of the plateau at $d=4.1\ln10=9.4$ mm.

\exbox{}{Why is the ratio of sums, and not the mean of the ratios, used when pooling runs?}\sol  The ratio of sums weights each run by its irrigation and gives the share of the total irrigation; a mean of ratios would give a run with almost no irrigation the same weight as a large one.

\exbox{}{The text says Equation~(\ref{eq:dose}) is phenomenological. Name two properties it was chosen for.}\sol  The correct limits (zero dose drains nothing; large dose passes on a constant share) and a single scale.
\banner{Hydraulic memory}\label{ch:mem}
\section{The idea}
A soil that starts a season wet releases some of its stored water during the season. Where does the released water go: to transpiration, to evaporation, or to drainage below the roots? Hydraulic memory is the share that drains. It is a property of the soil, not of one season.

\section{Definition: memory fate as a pooled slope}
Experiment E9 imposes, for each soil, an initial suction $h_0$ (ten values from 30 to 3,000 cm) and a rainfall scale $P$ (ten values from 0.5 to 1.4): a $10\times10$ lattice, 100 runs per soil. Wetter initial states (small $h_0$) release more storage, so the released storage $\bar R$ varies along the $h_0$ axis. For a flux $x\in\{E_s,T,D,\mathrm{oth}\}$ the memory fate is the within-$P$ pooled regression slope of that flux on $\bar R$:
\begin{equation}\label{eq:phi}
\boxed{\;\phi_x=\frac{\sum_P\sum_{c\in P}(\bar R_c-\bar R_P)(x_c-\bar x_P)}{\sum_P\sum_{c\in P}(\bar R_c-\bar R_P)^2},\qquad\phi_{E_s}+\phi_T+\phi_D+\phi_{\mathrm{oth}}=1\;}
\end{equation}
\par\noindent where $\phi_x$ is the fraction of the variation of rainfall that goes to the flux $x$ ($x=E_s,T,D$ or the rest), the sums run over the patterns $P$ and the cells $c$ of each, $\bar R_c$ is the mean rainfall of cell $c$, $\bar R_P$ the mean of pattern $P$, $x_c$ the value of the flux in the cell and $\bar x_P$ its mean over the pattern.

The sums run over lattice cells $c$ of one soil, and bars denote means within each rain level $P$. Pooling within $P$ removes the first-order effect of the change of rainfall on $\bar R$ and isolates the response to the antecedent state. $\phi_D$ is the hydraulic memory of drainage: the fraction of an additional millimetre of antecedent storage released that is shed as drainage. The unity constraint follows from the water balance when the irrigation response to $\bar R$ is negligible within the lattice. Confidence intervals come from bootstrapping the 100 cells (1,000 resamples).

\exbox{A slope by hand}{Within one rain level take four cells with released storage $\bar R=(10,20,30,40)$ mm and drainage $D=(6,13,17,24)$ mm. Means: $\bar R_P=25$, $\bar D_P=15$. Deviations of $\bar R$: $(-15,-5,5,15)$; of $D$: $(-9,-2,2,9)$. The numerator is $135+10+10+135=290$; the denominator is $225+25+25+225=500$. So $\phi_D=0.58$: of each extra millimetre released, 0.58 mm drains. In a real analysis the sums run over all ten rain levels of the lattice.}



\fig{p1_mem}{0.9}{Hydraulic memory of drainage across the ten soils, from the ensemble.}

\begin{kfig}{Hydraulic memory: the response to a saving decays as the store drains; the memory fate $\phi$ is the share of the saving still held by the soil after a season.}
\draw[axis] (0,0)--(7.4,0);\draw[axis] (0,0)--(0,4.2);
\node[lab,below] at (3.7,-0.45) {Time after the saving};\node[lab,rotate=90] at (-0.8,2.1) {Retained saving};
\draw[crv,domain=0:7,samples=40,smooth] plot (\x,{3.8*exp(-\x/5)});
\draw[ln,domain=0:7,samples=40,smooth] plot (\x,{3.8*exp(-\x/1.6)});
\draw[gd] (6.0,0)--(6.0,{3.8*exp(-6/5)});\node[lab,above] at (6.0,{3.8*exp(-6/5)+0.05}) {$\phi$};
\node[lab,text=khmag] at (5.2,2.5) {fine soil};\node[lab,text=khcyan!70!black] at (2.1,1.0) {coarse soil};
\end{kfig}

\section{Results across soils}
The memory fate $\phi_D$ runs from 0.85 in sand (S03) to 0.02 in silty clay loam (S10); see Table~\ref{tab:soils}. It follows the pore-size parameter $n$ (Spearman correlation $0.95$) and the saturated conductivity ($0.97$), and it transfers to independent evaporation and mulch experiments (correlations $0.85$ and $0.89$).

\section{The rate-competition law}
Let stored water $S$ above field capacity leave by two parallel first-order sinks: drainage with rate constant $k_D\propto\Ks/L$ and withdrawal by evaporation and uptake with $k_W\propto q^*/L$, where $L$ is a common length. Then $dS/dt=-(k_D+k_W)S$ and the drainage that leaves over the whole release is
\begin{equation}\label{eq:comp}\boxed{\int_0^\infty k_DS\,dt=\frac{k_D}{k_D+k_W}S_0\quad\Longrightarrow\quad\boxed{\;\phi_D=\frac{\Ks}{\Ks+q^*}\;}}\end{equation}
\par\noindent where $k_D$ and $k_W$ are the rate constants of drainage and of the other withdrawals (d$^{-1}$), $S$ the storage that remains at time $t$ (mm), $S_0$ the initial storage, $\phi_D$ the fraction that drains, $K_s$ the saturated conductivity (mm\,d$^{-1}$) and $q^*$ the mean withdrawal rate of the crop and soil (mm\,d$^{-1}$).

The law is bounded in $[0,1]$, tends to 0 for an impermeable soil and to 1 for a free-draining one, and is the only one-parameter form with these limits that is linear in $\Ks$ at small $\Ks$. By the Buckingham theorem $\phi_D=F(\Ks/q^*,n,\alpha h^*)$, and the law keeps the first group only. The constant $q^*$ is one withdrawal scale shared by all soils, fitted by nonlinear least squares on the eight soils with a literature $\Ks$. Left-out soils are predicted with a leave-one-soil-out root-mean-square error of 0.058.

\exbox{Reading the law}{If $\Ks=q^*$ then $\phi_D=1/2$: drainage and withdrawal are equally fast. A soil with $\Ks=9q^*$ has $\phi_D=0.9$, and one with $\Ks=q^*/9$ has $\phi_D=0.1$. A ninefold step in conductivity moves the memory fate from 0.1 to 0.5 to 0.9, which is why the memory of sand and silty clay loam are so far apart.}



\section{Why the memory transfers to other experiments}
The same competition decides where any extra water goes, whether it comes from initial storage, from a lower evaporation factor or from a mulch; the share is a property of the conductivity of the soil, not of the source of the water.

\section{Climate and the recoverable share of a saving}
In experiment E16, a $2^6$ factorial per soil (rainfall scale, temperature shift, trigger, dose, nitrogen rate, soil-evaporation factor), paired differences between a high and a low evaporation factor give the loss fate of each component. The results: warming raises the evaporative demand, so withdrawal is faster and less of a saved millimetre drains (the recoverable share rises); a wetter climate keeps the profile closer to field capacity, where drainage is fast, so more drains. The nitrogen rate changes the growth of the crop but not the hydraulic properties; it has no effect and so serves as a \emph{negative control}.

\begin{daisybox}{how the tests were done}
Leave-one-soil-out prediction compares four soil predictors (total available water, $\Se$ at 100 cm, $n$ and $\log\Ks$) and a mean-only benchmark. Spearman correlations come with exact permutation $p$ values and 10,000 bootstrap resamples; an exact sign test across soils checks the direction. All inputs are annual values read from the Daisy logs of E9, E7 and E16.
\end{daisybox}

\exbox{}{Four cells have $\bar R=(5,15,25,35)$ and $E_s=(20,22,24,26)$ in one rain level. Compute $\phi_{E_s}$.}\sol  Deviations $(-15,-5,5,15)$ and $(-3,-1,1,3)$: numerator $45+5+5+45=100$, denominator $500$, $\phi_{E_s}=0.20$.

\exbox{}{Derive Equation~(\ref{eq:comp}) from the two first-order sinks, showing the integral.}\sol  $dS/dt=-(k_D+k_W)S\Rightarrow S=S_0e^{-(k_D+k_W)t}$; $\int_0^\infty k_DS\,dt=k_DS_0/(k_D+k_W)$; with $k_D\propto\Ks$ and $k_W\propto q^*$ the factor $L$ cancels.

\exbox{}{If $q^*=1.5$ cm\,d$^{-1}$ (a hypothetical value) find $\phi_D$ for $\Ks=0.5$, $5$ and $50$ cm\,d$^{-1}$.}\sol  $\phi_D=0.25,\ 0.77,\ 0.97$ for $\Ks=0.5,\ 5,\ 50$ (with $q^*=1.5$).

\exbox{}{Explain why a nitrogen factor with no effect is useful as a negative control.}\sol  An effect that is known to be absent shows what a null result looks like and guards against reading a pattern into noise.

\exbox{}{A mulch saves 40 mm of evaporation in sand ($\phi_D=0.85$) and in silty clay loam ($\phi_D=0.02$). Estimate, using the memory fate as the drained share, how many millimetres stay available to the crop in each.}\sol  Sand: $40\times(1-0.85)=6$ mm stay; silty clay loam: $40\times(1-0.02)=39.2$ mm stay.
\banner{Irrigation control layers: sensor, trigger and trigger law}\label{ch:ctrl}
\section{Sensor and trigger}
A simulated tensiometer records the suction $h$ at the sensor depth every day. For year $y$ the layer counts the growing-season days (183 from 1 April) beyond the trigger $h_T$:
\begin{equation}\boxed{N_y(h_T)=\#\{d\in\text{season}_y:h_d>h_T\},\qquad X=100\Bigl(1-\frac{\bar N_{\rm irr}}{\bar N_{\rm rainfed}}\Bigr),\ \bar N=\frac1{20}\sum_yN_y.}\end{equation}
\par\noindent where $N_y(h_T)$ is the number of days of year $y$ on which the suction $h_d$ exceeds the trigger $h_T$ (cm), $\bar N_{\rm irr}$ and $\bar N_{\rm rainfed}$ the means over the 20 years for the irrigated and the rainfed runs, and $X$ the percentage reduction of the days beyond the trigger.


\exbox{Days beyond the trigger}{The layer prints: ``Irrigated: 85 of 183 days beyond 300 cm. Rainfed: 168 of 183 days beyond. Irrigation removes 39 \%.'' The 39 per cent is a 20-year mean while 85 and 168 are for the year drawn. (a) Verify the arithmetic for the year drawn. (b) If the rainfed 20-year mean is $\bar N_{\rm rainfed}=168$, what is the irrigated mean?}\sol  (a) For the year drawn, $X=100(1-85/168)=49.4$ per cent, which differs from the printed 39 per cent because $X$ uses the 20-year means, not one year. (b) $39=100(1-\bar N_{\rm irr}/168)$ gives $\bar N_{\rm irr}=168\times0.61=102.5$ days. So in an average year the irrigated sensor is beyond the trigger on 102 of 183 days against 168 rainfed. The year drawn is wetter than average for the irrigated trace by $102.5-85=17.5$ days.

\subhd{Reading.} A trigger of $-300$ cm does not mean ``never beyond 300 cm'': irrigation starts \emph{after} the sensor crosses the trigger, the dose is applied once per event, and the suction rises again between events.


\exbox{The delta method}{When the settings differ from the stored ones, the record is moved by the delta method
\begin{equation}\boxed{h^{\rm x}_d=h_{\rm soil}\Bigl(f_{\rm ref}(h^{\rm rec}_d)+\frac17\sum_{q=-3}^{3}\bigl[f^{\rm set}_{d+q}-f^{\rm std}_{d+q}\bigr]\Bigr).}\end{equation}
\par\noindent where $h^{\rm x}_d$ is the suction predicted for day $d$ (cm), $h_{\rm soil}$ the scale of the soil (cm), $f_{\rm ref}$ the reference depletion curve evaluated at the recorded suction $h^{\rm rec}_d$, and $f^{\rm set}$, $f^{\rm std}$ the depletion fractions of the chosen set and of the standard set, averaged over the seven days $d-3,\ldots,d+3$.

Suppose on day $d$ the stored record has $h^{\rm rec}=100$ cm in the reference soil (Coarse loam, depletion $f_{\rm ref}=f(100)$ in that soil) and the target soil is the loam S06. The seven-day differences $f^{\rm set}-f^{\rm std}$ of the sensor layer (the daily model run at your settings minus at standard settings) are $0.02,0.03,0.05,0.06,0.04,0.03,0.02$. Compute the moved suction. (Assume that the stored record is at $f_{\rm ref}=0.30$.)}\sol  The mean difference is $(0.02+0.03+0.05+0.06+0.04+0.03+0.02)/7=0.25/7=0.0357$. The new depletion is $0.30+0.0357=0.3357$. Convert it to a suction in the loam by inverting Equation~(\ref{eq:fdep}) by bisection: $h^{\rm x}=208$ cm. The same depletion $0.30$ would be $h=190$ cm in the loam. The settings therefore raise the suction at the sensor from 190 to 208 cm, an increase of 9 per cent for a depletion change of 0.036. Because suction is steep in depletion, a small change in depletion is a large change in suction.


\begin{kfig}{Relative yield against trigger depletion: yield stays near its maximum for small depletions and then falls; the largest trigger with $Y_{rel}\ge0.95$ is read from the dashed lines.}
\draw[axis] (0,0)--(7.4,0);\draw[axis] (0,0)--(0,4.4);
\node[lab,below] at (3.7,-0.5) {Trigger depletion $f$};\node[lab,rotate=90] at (-0.9,2.2) {Relative yield $Y_{rel}$};
\foreach \x/\t in {0/0,2.4/0.25,4.8/0.5,7.2/0.75}{\draw[ln] (\x,0)--(\x,-0.1);\node[lab,below] at (\x,-0.1) {\t};}
\draw[crv,smooth] plot coordinates {(0,4.0)(1.2,3.98)(2.4,3.8)(3.6,3.4)(4.8,2.8)(6,2.1)(7.2,1.3)};
\draw[gd] (0,3.8)--(7.2,3.8);\node[lab,left] at (0,3.8) {0.95};\draw[gd] (2.4,0)--(2.4,3.8);
\end{kfig}

\section{Trigger law}
The suction at which a soil reaches a chosen depletion is $h(f)=\{h:f(h)=f\}$ (Equation~\ref{eq:fdep} inverted), the net irrigation per event is $Z_n=f\cdot\mathrm{AWC}\cdot R_{\max}$, and the relative yield at a trigger is
\begin{equation}\boxed{Y_{\rm rel}(h_T)=\frac{Y(h_T)}{\max_hY(h)}\quad(\text{data-driven, run set E2}).}\end{equation}
\par\noindent where $Y_{\rm rel}$ is the relative yield (0 to 1), $Y(h_T)$ the yield for trigger $h_T$ (Mg\,ha$^{-1}$) and $\max_hY(h)$ its maximum over all triggers, taken from the run set E2.


\exbox{Relative yield against trigger depletion}{Use the archived E2 trigger runs to compute, for three soils, the depletion at each trigger and the relative yield, and find the largest trigger with $Y_{\rm rel}\ge0.95$.}\sol  For each soil take $Y_{\max}$ as the largest yield over its triggers, and evaluate $f(h_T)$ with Equation~(\ref{eq:fdep}). Table~\ref{tab:yrel} lists the result.\begin{table}[!htbp]\centering\small\caption{Archived trigger runs (E2): depletion at the trigger, irrigation, yield and relative yield.}\label{tab:yrel}
\begin{tabular}{lrrrrr}\toprule
\hd{Soil} & \hd{Trigger (cm)} & \hd{$f$} & \hd{Irrigation (mm)} & \hd{Yield (Mg/ha)} & \hd{$Y_{\rm rel}$}\\\midrule
S06 & 50 & 0.000 & 424 & 6.62 & 1.000\\
S06 & 100 & 0.000 & 346 & 6.61 & 0.997\\
S06 & 200 & 0.321 & 257 & 6.38 & 0.964\\
S06 & 300 & 0.468 & 212 & 6.12 & 0.925\\
S06 & 500 & 0.614 & 173 & 5.96 & 0.899\\
S06 & 1000 & 0.759 & 116 & 5.19 & 0.783\\
S06 & 3000 & 0.899 & 37 & 3.57 & 0.540\\
S03 & 50 & 0.000 & 122 & 3.04 & 1.000\\
S03 & 100 & 0.000 & 6 & 1.23 & 0.406\\
S03 & 200 & 0.688 & 0 & 1.13 & 0.373\\
S03 & 300 & 0.842 & 0 & 1.13 & 0.373\\
S03 & 500 & 0.933 & 0 & 1.13 & 0.373\\
S03 & 1000 & 0.979 & 0 & 1.13 & 0.373\\
S03 & 3000 & 0.997 & 0 & 1.13 & 0.373\\
S10 & 50 & 0.000 & 383 & 5.99 & 0.908\\
S10 & 100 & 0.000 & 347 & 6.23 & 0.945\\
S10 & 200 & 0.142 & 306 & 6.51 & 0.988\\
S10 & 300 & 0.237 & 278 & 6.58 & 0.998\\
S10 & 500 & 0.361 & 242 & 6.51 & 0.987\\
S10 & 1000 & 0.522 & 190 & 6.42 & 0.975\\
S10 & 3000 & 0.745 & 150 & 5.87 & 0.891\\\bottomrule\end{tabular}\end{table}

Loam: the largest trigger with $Y_{\rm rel}\ge0.95$ is 200 cm, $f=0.32$. Sand: the largest trigger with $Y_{\rm rel}\ge0.95$ is 50 cm, $f=0.00$. Silty clay loam: the largest trigger with $Y_{\rm rel}\ge0.95$ is 1500 cm, $f=0.61$. 

\subhd{Reading.} In suction the thresholds differ by a large factor between soils; in depletion they cluster near $f^*\approx0.45$ (the study's soil-independent value; 5 per cent yield loss), which is the content of the trigger law. The statement is data-driven for one crop, one weather series and a sensor at 20 cm.


\exbox{How fast each soil dries}{Every soil is compared on the same crop demand $ET_c$; only the water held differs:
\begin{equation}\boxed{r_s=\frac{ET_c}{\mathrm{TAW}_s},\quad\mathrm{TAW}_s=\mathrm{AWC}_sz_r,\quad t_{{\rm trig},s}=\frac{f_{{\rm trig},s}}{r_s}.}\end{equation}
\par\noindent where $r_s$ is the depletion rate of soil $s$ (fraction of TAW per day), $ET_c$ the crop evapotranspiration (mm\,d$^{-1}$), ${\rm TAW}_s$ the total available water of the soil (mm), ${\rm AWC}_s$ its available water capacity (mm\,cm$^{-1}$), $z_r$ the rooting depth (cm), $f_{{\rm trig},s}$ the depletion fraction at which the trigger fires and $t_{{\rm trig},s}$ the time to the trigger (d).

Take $ET_c=3.5$ mm d$^{-1}$ (assumed peak-season value), $z_r=90$ cm and $f_{\rm trig}=0.45$. Compute $r_s$ in per cent per day and the days to the trigger for the ten soils.}\sol  For the sand, $\mathrm{TAW}=0.6\times90=54$ mm, $r=3.5/54=0.0648$, that is $6.5$ per cent per day, and $t=0.45/0.0648=6.9$ days. For the coarse loam, $\mathrm{TAW}=1.5\times90=135$ mm, $r=2.6$ per cent per day and $t=17.4$ days. These agree with the application text (about 7 days for sand and 18 days for coarse loam). Table~\ref{tab:dry} gives all ten soils.\begin{table}[!htbp]\centering\small\caption{Speed of depletion and days from field capacity to the trigger $f=0.45$ (assumed $ET_c=3.5$ mm d$^{-1}$, $z_r=90$ cm).}\label{tab:dry}
\begin{tabular}{lrrrrr}\toprule
\hd{ID} & \hd{Soil} & \hd{AWC} & \hd{TAW (mm)} & \hd{$100r_s$ (\%/d)} & \hd{Days to trigger}\\\midrule
S01 & Coarse loam & 1.5 & 135 & 2.6 & 17.4\\
S02 & Coarse sand & 0.5 & 45 & 7.8 & 5.8\\
S03 & Sand & 0.6 & 54 & 6.5 & 6.9\\
S04 & Loamy sand & 0.9 & 81 & 4.3 & 10.4\\
S05 & Sandy loam & 1.2 & 108 & 3.2 & 13.9\\
S06 & Loam & 1.7 & 153 & 2.3 & 19.7\\
S07 & Silt loam & 2.0 & 180 & 1.9 & 23.1\\
S08 & Sandy clay loam & 1.4 & 126 & 2.8 & 16.2\\
S09 & Clay loam & 1.6 & 144 & 2.4 & 18.5\\
S10 & Silty clay loam & 1.7 & 153 & 2.3 & 19.7\\\bottomrule\end{tabular}\end{table}
\textbf{Reading.} The coarse sand dries fastest (smallest AWC, 0.5 mm cm$^{-1}$) and the silt loam slowest (AWC 2.0): a four-fold spread in the interval between irrigations, from 5.8 to 23 days. A sprinkler round that takes three days is long compared with the 5.8 days of the coarse sand but short compared with the 23 days of the silt loam, which is why the f=0.45 advice for coarse soils exists.


\exbox{Choosing between f = 0.50 and f = 0.45}{The application advises $f=0.45$ if any of the following holds: a round of the design takes more than 2 days; the soil is sandy (sand $\ge60$ per cent or AWC $\le0.9$ mm cm$^{-1}$); the crop is at $\ge85$ per cent cover (critical stage). Otherwise $f=0.50$. Decide for three cases.}
\textbf{Case A:} maize on the loam (sand 40 per cent, AWC 1.7), drip system, round 1.5 days, date with cover 60 per cent. No $f=0.45$ condition holds (round $\le2$ d, sand $<60$, AWC $>0.9$, cover $<85$ per cent), so the advice is $f=0.50$. In suction, $f=0.50$ in the loam is $h=332$ cm.

\subhd{Case B:} same field but a sprinkler round of 3.5 days. The round exceeds 2 days, so the advice is $f=0.45$ ($h=284$ cm): the field must be triggered earlier because the last block is irrigated 3.5 days after the first, and during that time depletion advances by $3.5\times r_s$. With $r_s=3.5/153=2.3$ per cent per day (loam, $\mathrm{TAW}=1.7\times90=153$ mm) the last block gains $0.08$ in depletion, so starting at $0.45$ the last block reaches $0.53$, just above $0.50$.

\subhd{Case C:} sand (92 per cent sand), drip, round 1 day. The soil is sandy, so $f=0.45$ ($h=143$ cm).


\exbox{Reading the sensor tile}{The Dashboard tile ``Sensor / trigger'' reads $0.02/0.38$ for the coarse loam on a date. Interpret it and find the suction of the sensor and the days to irrigation at $r=2.6$ per cent per day.}\sol  The first number is the present depletion of the sensor layer ($f=0.02$, almost field capacity); the second is the trigger depletion $f_T=f(300\text{ cm})=0.38$ of the coarse loam. Irrigation is due when the first reaches the second. The sensor suction is $h(0.02)=105 cm$, close to the field capacity of 100 cm. At $r=0.026$ d$^{-1}$ the soil needs $(0.38-0.02)/0.026=13.8$ days without rain to reach the trigger.
\IfFileExists{extra/long03.tex}{\subsection{Complete calculation of a trigger suction}
\exbox{Trigger suction by bisection}{Find the suction at which the loam S06 has lost 50 per cent of its readily available water, showing the first eight bisection steps, and compare three soils at three trigger fractions.}\sol
\ans The trigger is the suction at which the readily available water has been used up; for the soil it is found by bisection on $\log_{10}h$ of the depletion function $f(h)=\dfrac{S_e(100)-S_e(h)}{S_e(100)-S_e(15849)}$ until it equals the trigger fraction
\ans Step 1: $\log_{10}h=2.3000$, $h=200$ cm, $f=0.3201$ $<$ 0.5
\ans Step 2: $\log_{10}h=3.4500$, $h=2818$ cm, $f=0.8935$ $\ge$ 0.5
\ans Step 3: $\log_{10}h=2.8750$, $h=750$ cm, $f=0.7053$ $\ge$ 0.5
\ans Step 4: $\log_{10}h=2.5875$, $h=387$ cm, $f=0.5456$ $\ge$ 0.5
\ans Step 5: $\log_{10}h=2.4437$, $h=278$ cm, $f=0.4422$ $<$ 0.5
\ans Step 6: $\log_{10}h=2.5156$, $h=328$ cm, $f=0.4961$ $<$ 0.5
\ans Step 7: $\log_{10}h=2.5516$, $h=356$ cm, $f=0.5214$ $\ge$ 0.5
\ans Step 8: $\log_{10}h=2.5336$, $h=342$ cm, $f=0.5089$ $\ge$ 0.5
\ans After 60 halvings the trigger suction of the loam at 50 per cent depletion is $h_{50}=332$ cm
\ans Soil S02 (Coarse sand): $h_{50}=303$ cm; at 30 per cent $179$ cm; at 70 per cent $646$ cm
\ans Soil S06 (Loam): $h_{50}=332$ cm; at 30 per cent $190$ cm; at 70 per cent $731$ cm
\ans Soil S10 (Silty clay loam): $h_{50}=906$ cm; at 30 per cent $389$ cm; at 70 per cent $2354$ cm \finans


}{}
\IfFileExists{extra/proj03.tex}{\banner{Design project 3: choosing a suction trigger for six soils}
\textbf{Brief.} A sensor-controlled system is to be installed on six soils of a farm. The crop (a vegetable) roots to 0.6 m and tolerates a depletion of 50 per cent of the available water before it is stressed; its peak demand is 5 mm\,d$^{-1}$. The designer must set the trigger of each sensor so that (i) the depletion stays below 50 per cent, (ii) the trigger is within the reach of the soil (taken as 80 per cent of $h_c$ to allow a margin), and (iii) the net dose $Z_n$ and the interval at the peak are known.

\textbf{Method.} For each soil: the suction at 50 per cent depletion $h_{50}$ is found by bisection of $f(h)=0.5$ with the van Genuchten curve; the reach $h_c$ from Equation~(\ref{eq:reach3}); the trigger is $h_t=\min(h_{50},0.8h_c)$; the actual depletion $f_t=f(h_t)$; the net dose $Z_n=f_t\cdot\mathrm{AWC}\cdot R_{max}$ with $R_{max}=60$ cm; and the interval at the peak $T=Z_n/5$.

\begin{table}[!htbp]\centering\small\caption{Design of the trigger for six soils: suction at 50 \% depletion, reach, chosen trigger, depletion at the trigger, net dose and interval at the peak demand.}
\begin{tabular}{llrrrrrr}\toprule
\hd{Soil} & \hd{Class} & \hd{$h_{50}$} & \hd{$h_c$} & \hd{$h_t$} & \hd{$f_t$} & \hd{$Z_n$ (mm)} & \hd{Interval (d)}\\\midrule
S03 & Sand & 151 & 199 & 151 & 0.50 & 18 & 3.6\\
S05 & Sandy loam & 217 & 1\,615 & 217 & 0.50 & 36 & 7.2\\
S06 & Loam & 332 & 9\,186 & 332 & 0.50 & 51 & 10.2\\
S07 & Silt loam & 483 & 29\,528 & 483 & 0.50 & 60 & 12.0\\
S09 & Clay loam & 618 & 47\,608 & 618 & 0.50 & 48 & 9.6\\
S10 & Silty clay loam & 906 & 127\,830 & 906 & 0.50 & 51 & 10.2\\\bottomrule\end{tabular}\end{table}

\textbf{Reading the table.} In coarse soils (sand, sandy loam) the depletion limit $h_{50}$ is low and is the binding constraint; the trigger is $h_{50}$ and the dose is small, so the sensors call for frequent small irrigations (the interval is a day or two). In the fine soils the reach is the binding limit when $h_{50}$ exceeds $0.8h_c$: the trigger cannot be set at a depletion of 50 per cent because the roots could not dry the soil that far, and the practical depletion at the trigger is lower than the design limit. The designer should then choose a drip or a pulsed system in the coarse soils and accept that the fine soils will be irrigated at less than their allowed depletion.

\textbf{Check.} For the loam, the trigger of 332 cm gives a depletion of 50 per cent; a net dose of 51 mm every 10.2 days at the peak. The drained share of such a dose is read from the dose law of this chapter.

}{}
\IfFileExists{extra/ten03.tex}{\subsection{Worked solutions on the ten soils}
The reach of the trigger decides whether irrigation can be scheduled by suction. The calculation of Eq.~(\ref{eq:reach}) is repeated for each soil.

\exbox{Ten-soil example 3a: Coarse loam (S01)}{Coarse loam: $\alpha=0.0961$ cm$^{-1}$, $n=1.353$, $\ell=0.5$, $\kappa=3.5\times10^{-10}$. Compute the exponent $E$, the reach $h_c$ and the product $\alpha h_c$, and say whether a trigger at $h_T=300$ cm can fire.}\sol \givenline{$\alpha=0.0961$ cm$^{-1}$ ; $n=1.353$ ; $\ell=0.5$ ; $\kappa=3.5\times10^{-10}$ ; $h_T=300$ cm can fire}\ans $m=0.261$; \par\indent We know that
\[E=2n+\ell(n-1)\ti{i}\]
Substituting the values in equation \rfi{i},
\[E=2(1.353)+0.5(0.353)=2.882\] \ans \[m^2/\kappa=0.068/(3.5\times10^{-10})=1.94\times10^{8}\] \ans \[h_c=\frac{1}{0.0961}(1.94\times10^{8})^{1/2.882}=7814\ \text{cm}\] \ans $\alpha h_c=751$ \ans A trigger of 300 cm is below the reach of 7814 cm: it can fire and irrigation will follow \finans

\exbox{Ten-soil example 3a: Coarse sand (S02)}{Coarse sand: $\alpha=0.0536$ cm$^{-1}$, $n=1.600$, $\ell=0.5$, $\kappa=3.5\times10^{-10}$. Compute the exponent $E$, the reach $h_c$ and the product $\alpha h_c$, and say whether a trigger at $h_T=300$ cm can fire.}\sol \givenline{$\alpha=0.0536$ cm$^{-1}$ ; $n=1.600$ ; $\ell=0.5$ ; $\kappa=3.5\times10^{-10}$ ; $h_T=300$ cm can fire}\ans $m=0.375$; \par\indent We know that
\[E=2n+\ell(n-1)\ti{i}\]
Substituting the values in equation \rfi{i},
\[E=2(1.600)+0.5(0.600)=3.500\] \ans \[m^2/\kappa=0.141/(3.5\times10^{-10})=4.02\times10^{8}\] \ans \[h_c=\frac{1}{0.0536}(4.02\times10^{8})^{1/3.500}=5359\ \text{cm}\] \ans $\alpha h_c=287$ \ans A trigger of 300 cm is below the reach of 5359 cm: it can fire and irrigation will follow \finans

\exbox{Ten-soil example 3a: Sand (S03)}{Sand: $\alpha=0.1450$ cm$^{-1}$, $n=2.680$, $\ell=0.5$, $\kappa=3.5\times10^{-10}$. Compute the exponent $E$, the reach $h_c$ and the product $\alpha h_c$, and say whether a trigger at $h_T=300$ cm can fire.}\sol \givenline{$\alpha=0.1450$ cm$^{-1}$ ; $n=2.680$ ; $\ell=0.5$ ; $\kappa=3.5\times10^{-10}$ ; $h_T=300$ cm can fire}\ans $m=0.627$; \par\indent We know that
\[E=2n+\ell(n-1)\ti{i}\]
Substituting the values in equation \rfi{i},
\[E=2(2.680)+0.5(1.680)=6.200\] \ans \[m^2/\kappa=0.393/(3.5\times10^{-10})=1.12\times10^{9}\] \ans \[h_c=\frac{1}{0.1450}(1.12\times10^{9})^{1/6.200}=199\ \text{cm}\] \ans $\alpha h_c=29$ \ans A trigger of 300 cm is beyond the reach of 199 cm: the soil cannot supply the demand before it, so the crop stresses first \finans

\exbox{Ten-soil example 3a: Loamy sand (S04)}{Loamy sand: $\alpha=0.1240$ cm$^{-1}$, $n=2.280$, $\ell=0.5$, $\kappa=3.5\times10^{-10}$. Compute the exponent $E$, the reach $h_c$ and the product $\alpha h_c$, and say whether a trigger at $h_T=300$ cm can fire.}\sol \givenline{$\alpha=0.1240$ cm$^{-1}$ ; $n=2.280$ ; $\ell=0.5$ ; $\kappa=3.5\times10^{-10}$ ; $h_T=300$ cm can fire}\ans $m=0.561$; \par\indent We know that
\[E=2n+\ell(n-1)\ti{i}\]
Substituting the values in equation \rfi{i},
\[E=2(2.280)+0.5(1.280)=5.200\] \ans \[m^2/\kappa=0.315/(3.5\times10^{-10})=9\times10^{8}\] \ans \[h_c=\frac{1}{0.1240}(9\times10^{8})^{1/5.200}=425\ \text{cm}\] \ans $\alpha h_c=53$ \ans A trigger of 300 cm is below the reach of 425 cm: it can fire and irrigation will follow \finans

\exbox{Ten-soil example 3a: Sandy loam (S05)}{Sandy loam: $\alpha=0.0750$ cm$^{-1}$, $n=1.890$, $\ell=0.5$, $\kappa=3.5\times10^{-10}$. Compute the exponent $E$, the reach $h_c$ and the product $\alpha h_c$, and say whether a trigger at $h_T=300$ cm can fire.}\sol \givenline{$\alpha=0.0750$ cm$^{-1}$ ; $n=1.890$ ; $\ell=0.5$ ; $\kappa=3.5\times10^{-10}$ ; $h_T=300$ cm can fire}\ans $m=0.471$; \par\indent We know that
\[E=2n+\ell(n-1)\ti{i}\]
Substituting the values in equation \rfi{i},
\[E=2(1.890)+0.5(0.890)=4.225\] \ans \[m^2/\kappa=0.222/(3.5\times10^{-10})=6.34\times10^{8}\] \ans \[h_c=\frac{1}{0.0750}(6.34\times10^{8})^{1/4.225}=1615\ \text{cm}\] \ans $\alpha h_c=121$ \ans A trigger of 300 cm is below the reach of 1615 cm: it can fire and irrigation will follow \finans

\exbox{Ten-soil example 3a: Loam (S06)}{Loam: $\alpha=0.0360$ cm$^{-1}$, $n=1.560$, $\ell=0.5$, $\kappa=3.5\times10^{-10}$. Compute the exponent $E$, the reach $h_c$ and the product $\alpha h_c$, and say whether a trigger at $h_T=300$ cm can fire.}\sol \givenline{$\alpha=0.0360$ cm$^{-1}$ ; $n=1.560$ ; $\ell=0.5$ ; $\kappa=3.5\times10^{-10}$ ; $h_T=300$ cm can fire}\ans $m=0.359$; \par\indent We know that
\[E=2n+\ell(n-1)\ti{i}\]
Substituting the values in equation \rfi{i},
\[E=2(1.560)+0.5(0.560)=3.400\] \ans \[m^2/\kappa=0.129/(3.5\times10^{-10})=3.68\times10^{8}\] \ans \[h_c=\frac{1}{0.0360}(3.68\times10^{8})^{1/3.400}=9186\ \text{cm}\] \ans $\alpha h_c=331$ \ans A trigger of 300 cm is below the reach of 9186 cm: it can fire and irrigation will follow \finans

\exbox{Ten-soil example 3a: Silt loam (S07)}{Silt loam: $\alpha=0.0200$ cm$^{-1}$, $n=1.410$, $\ell=0.5$, $\kappa=3.5\times10^{-10}$. Compute the exponent $E$, the reach $h_c$ and the product $\alpha h_c$, and say whether a trigger at $h_T=300$ cm can fire.}\sol \givenline{$\alpha=0.0200$ cm$^{-1}$ ; $n=1.410$ ; $\ell=0.5$ ; $\kappa=3.5\times10^{-10}$ ; $h_T=300$ cm can fire}\ans $m=0.291$; \par\indent We know that
\[E=2n+\ell(n-1)\ti{i}\]
Substituting the values in equation \rfi{i},
\[E=2(1.410)+0.5(0.410)=3.025\] \ans \[m^2/\kappa=0.085/(3.5\times10^{-10})=2.42\times10^{8}\] \ans \[h_c=\frac{1}{0.0200}(2.42\times10^{8})^{1/3.025}=29528\ \text{cm}\] \ans $\alpha h_c=591$ \ans A trigger of 300 cm is below the reach of 29528 cm: it can fire and irrigation will follow \finans

\exbox{Ten-soil example 3a: Sandy clay loam (S08)}{Sandy clay loam: $\alpha=0.0590$ cm$^{-1}$, $n=1.480$, $\ell=0.5$, $\kappa=3.5\times10^{-10}$. Compute the exponent $E$, the reach $h_c$ and the product $\alpha h_c$, and say whether a trigger at $h_T=300$ cm can fire.}\sol \givenline{$\alpha=0.0590$ cm$^{-1}$ ; $n=1.480$ ; $\ell=0.5$ ; $\kappa=3.5\times10^{-10}$ ; $h_T=300$ cm can fire}\ans $m=0.324$; \par\indent We know that
\[E=2n+\ell(n-1)\ti{i}\]
Substituting the values in equation \rfi{i},
\[E=2(1.480)+0.5(0.480)=3.200\] \ans \[m^2/\kappa=0.105/(3.5\times10^{-10})=3.01\times10^{8}\] \ans \[h_c=\frac{1}{0.0590}(3.01\times10^{8})^{1/3.200}=7559\ \text{cm}\] \ans $\alpha h_c=446$ \ans A trigger of 300 cm is below the reach of 7559 cm: it can fire and irrigation will follow \finans

\exbox{Ten-soil example 3a: Clay loam (S09)}{Clay loam: $\alpha=0.0190$ cm$^{-1}$, $n=1.310$, $\ell=0.5$, $\kappa=3.5\times10^{-10}$. Compute the exponent $E$, the reach $h_c$ and the product $\alpha h_c$, and say whether a trigger at $h_T=300$ cm can fire.}\sol \givenline{$\alpha=0.0190$ cm$^{-1}$ ; $n=1.310$ ; $\ell=0.5$ ; $\kappa=3.5\times10^{-10}$ ; $h_T=300$ cm can fire}\ans $m=0.237$; \par\indent We know that
\[E=2n+\ell(n-1)\ti{i}\]
Substituting the values in equation \rfi{i},
\[E=2(1.310)+0.5(0.310)=2.775\] \ans \[m^2/\kappa=0.056/(3.5\times10^{-10})=1.6\times10^{8}\] \ans \[h_c=\frac{1}{0.0190}(1.6\times10^{8})^{1/2.775}=47608\ \text{cm}\] \ans $\alpha h_c=905$ \ans A trigger of 300 cm is below the reach of 47608 cm: it can fire and irrigation will follow \finans

\exbox{Ten-soil example 3a: Silty clay loam (S10)}{Silty clay loam: $\alpha=0.0100$ cm$^{-1}$, $n=1.230$, $\ell=0.5$, $\kappa=3.5\times10^{-10}$. Compute the exponent $E$, the reach $h_c$ and the product $\alpha h_c$, and say whether a trigger at $h_T=300$ cm can fire.}\sol \givenline{$\alpha=0.0100$ cm$^{-1}$ ; $n=1.230$ ; $\ell=0.5$ ; $\kappa=3.5\times10^{-10}$ ; $h_T=300$ cm can fire}\ans $m=0.187$; \par\indent We know that
\[E=2n+\ell(n-1)\ti{i}\]
Substituting the values in equation \rfi{i},
\[E=2(1.230)+0.5(0.230)=2.575\] \ans \[m^2/\kappa=0.035/(3.5\times10^{-10})=9.99\times10^{7}\] \ans \[h_c=\frac{1}{0.0100}(9.99\times10^{7})^{1/2.575}=127830\ \text{cm}\] \ans $\alpha h_c=1278$ \ans A trigger of 300 cm is below the reach of 127830 cm: it can fire and irrigation will follow \finans

}{}
\IfFileExists{extra/tenb03.tex}{\section{Worked solutions on ten crops}
The ten crops of the library differ in their demand, their roots and their sensitivity. The same type of question is answered for each of them at $ET_0=5.5$ mm\,d$^{-1}$; the data are typical values of the crops (Doorenbos--Kassam and FAO 56).

\exbox{Crop example 3a: Maize}{Maize on the loam S06: the trigger fires at the depletion fraction $p=0.55$. Using the demand $r=ET_c/TAW$ and $t_{\rm trig}=p/r$, find the time from full profile to the trigger and the number of irrigations in a season of 125 days.}\sol $TAW=1.7\times100=170$ mm; $r=1.20\times5.5/170=0.0388$ d$^{-1}$ \ans $t_{\rm trig}=0.55/0.0388=14.2$ d \ans Irrigations $\approx125/14.2=9$ \ans

\exbox{Crop example 3a: Wheat}{Wheat on the loam S06: the trigger fires at the depletion fraction $p=0.55$. Using the demand $r=ET_c/TAW$ and $t_{\rm trig}=p/r$, find the time from full profile to the trigger and the number of irrigations in a season of 150 days.}\sol $TAW=1.7\times120=204$ mm; $r=1.15\times5.5/204=0.0310$ d$^{-1}$ \ans $t_{\rm trig}=0.55/0.0310=17.7$ d \ans Irrigations $\approx150/17.7=8$ \ans

\exbox{Crop example 3a: Cotton}{Cotton on the loam S06: the trigger fires at the depletion fraction $p=0.65$. Using the demand $r=ET_c/TAW$ and $t_{\rm trig}=p/r$, find the time from full profile to the trigger and the number of irrigations in a season of 180 days.}\sol $TAW=1.7\times110=187$ mm; $r=1.15\times5.5/187=0.0338$ d$^{-1}$ \ans $t_{\rm trig}=0.65/0.0338=19.2$ d \ans Irrigations $\approx180/19.2=9$ \ans

\exbox{Crop example 3a: Tomato}{Tomato on the loam S06: the trigger fires at the depletion fraction $p=0.40$. Using the demand $r=ET_c/TAW$ and $t_{\rm trig}=p/r$, find the time from full profile to the trigger and the number of irrigations in a season of 135 days.}\sol $TAW=1.7\times70=119$ mm; $r=1.15\times5.5/119=0.0532$ d$^{-1}$ \ans $t_{\rm trig}=0.40/0.0532=7.5$ d \ans Irrigations $\approx135/7.5=18$ \ans

\exbox{Crop example 3a: Potato}{Potato on the loam S06: the trigger fires at the depletion fraction $p=0.35$. Using the demand $r=ET_c/TAW$ and $t_{\rm trig}=p/r$, find the time from full profile to the trigger and the number of irrigations in a season of 130 days.}\sol $TAW=1.7\times50=85$ mm; $r=1.15\times5.5/85=0.0744$ d$^{-1}$ \ans $t_{\rm trig}=0.35/0.0744=4.7$ d \ans Irrigations $\approx130/4.7=28$ \ans

\exbox{Crop example 3a: Soybean}{Soybean on the loam S06: the trigger fires at the depletion fraction $p=0.50$. Using the demand $r=ET_c/TAW$ and $t_{\rm trig}=p/r$, find the time from full profile to the trigger and the number of irrigations in a season of 130 days.}\sol $TAW=1.7\times60=102$ mm; $r=1.15\times5.5/102=0.0620$ d$^{-1}$ \ans $t_{\rm trig}=0.50/0.0620=8.1$ d \ans Irrigations $\approx130/8.1=16$ \ans

\exbox{Crop example 3a: Sugar beet}{Sugar beet on the loam S06: the trigger fires at the depletion fraction $p=0.55$. Using the demand $r=ET_c/TAW$ and $t_{\rm trig}=p/r$, find the time from full profile to the trigger and the number of irrigations in a season of 180 days.}\sol $TAW=1.7\times90=153$ mm; $r=1.20\times5.5/153=0.0431$ d$^{-1}$ \ans $t_{\rm trig}=0.55/0.0431=12.8$ d \ans Irrigations $\approx180/12.8=14$ \ans

\exbox{Crop example 3a: Sunflower}{Sunflower on the loam S06: the trigger fires at the depletion fraction $p=0.45$. Using the demand $r=ET_c/TAW$ and $t_{\rm trig}=p/r$, find the time from full profile to the trigger and the number of irrigations in a season of 125 days.}\sol $TAW=1.7\times80=136$ mm; $r=1.10\times5.5/136=0.0445$ d$^{-1}$ \ans $t_{\rm trig}=0.45/0.0445=10.1$ d \ans Irrigations $\approx125/10.1=12$ \ans

\exbox{Crop example 3a: Onion}{Onion on the loam S06: the trigger fires at the depletion fraction $p=0.30$. Using the demand $r=ET_c/TAW$ and $t_{\rm trig}=p/r$, find the time from full profile to the trigger and the number of irrigations in a season of 150 days.}\sol $TAW=1.7\times35=60$ mm; $r=1.05\times5.5/60=0.0971$ d$^{-1}$ \ans $t_{\rm trig}=0.30/0.0971=3.1$ d \ans Irrigations $\approx150/3.1=49$ \ans

\exbox{Crop example 3a: Alfalfa}{Alfalfa on the loam S06: the trigger fires at the depletion fraction $p=0.55$. Using the demand $r=ET_c/TAW$ and $t_{\rm trig}=p/r$, find the time from full profile to the trigger and the number of irrigations in a season of 100 days.}\sol $TAW=1.7\times100=170$ mm; $r=1.15\times5.5/170=0.0372$ d$^{-1}$ \ans $t_{\rm trig}=0.55/0.0372=14.8$ d \ans Irrigations $\approx100/14.8=7$ \ans

}{}
\IfFileExists{extra/teny03.tex}{\section{Worked solutions on the ten soils (third scenario)}
The same questions are repeated for different working conditions, to show how the answer moves with them.

\exbox{Ten-soil example 3c: Coarse loam (S01)}{Coarse loam: $\alpha=0.0961$ cm$^{-1}$, $n=1.353$, $\ell=0.5$, $\kappa=3.5\times10^{-10}$. Compute the exponent $E$, the reach $h_c$ and the product $\alpha h_c$, and say whether a trigger at $h_T=300$ cm can fire.}\sol $m=0.261$; $E=2n+\ell(n-1)=2(1.353)+0.5(0.353)=2.882$ \ans $m^2/\kappa=0.068/(3.5\times10^{-10})=1.94e+08$ \ans $h_c=\frac{1}{0.0961}(1.94e+08)^{1/2.882}=7814$ cm \ans $\alpha h_c=751$ \ans A trigger of 300 cm is below the reach of 7814 cm: it can fire and irrigation will follow \ans

\exbox{Ten-soil example 3c: Coarse sand (S02)}{Coarse sand: $\alpha=0.0536$ cm$^{-1}$, $n=1.600$, $\ell=0.5$, $\kappa=3.5\times10^{-10}$. Compute the exponent $E$, the reach $h_c$ and the product $\alpha h_c$, and say whether a trigger at $h_T=300$ cm can fire.}\sol $m=0.375$; $E=2n+\ell(n-1)=2(1.600)+0.5(0.600)=3.500$ \ans $m^2/\kappa=0.141/(3.5\times10^{-10})=4.02e+08$ \ans $h_c=\frac{1}{0.0536}(4.02e+08)^{1/3.500}=5359$ cm \ans $\alpha h_c=287$ \ans A trigger of 300 cm is below the reach of 5359 cm: it can fire and irrigation will follow \ans

\exbox{Ten-soil example 3c: Sand (S03)}{Sand: $\alpha=0.1450$ cm$^{-1}$, $n=2.680$, $\ell=0.5$, $\kappa=3.5\times10^{-10}$. Compute the exponent $E$, the reach $h_c$ and the product $\alpha h_c$, and say whether a trigger at $h_T=300$ cm can fire.}\sol $m=0.627$; $E=2n+\ell(n-1)=2(2.680)+0.5(1.680)=6.200$ \ans $m^2/\kappa=0.393/(3.5\times10^{-10})=1.12e+09$ \ans $h_c=\frac{1}{0.1450}(1.12e+09)^{1/6.200}=199$ cm \ans $\alpha h_c=29$ \ans A trigger of 300 cm is beyond the reach of 199 cm: the soil cannot supply the demand before it, so the crop stresses first \ans

\exbox{Ten-soil example 3c: Loamy sand (S04)}{Loamy sand: $\alpha=0.1240$ cm$^{-1}$, $n=2.280$, $\ell=0.5$, $\kappa=3.5\times10^{-10}$. Compute the exponent $E$, the reach $h_c$ and the product $\alpha h_c$, and say whether a trigger at $h_T=300$ cm can fire.}\sol $m=0.561$; $E=2n+\ell(n-1)=2(2.280)+0.5(1.280)=5.200$ \ans $m^2/\kappa=0.315/(3.5\times10^{-10})=9e+08$ \ans $h_c=\frac{1}{0.1240}(9e+08)^{1/5.200}=425$ cm \ans $\alpha h_c=53$ \ans A trigger of 300 cm is below the reach of 425 cm: it can fire and irrigation will follow \ans

\exbox{Ten-soil example 3c: Sandy loam (S05)}{Sandy loam: $\alpha=0.0750$ cm$^{-1}$, $n=1.890$, $\ell=0.5$, $\kappa=3.5\times10^{-10}$. Compute the exponent $E$, the reach $h_c$ and the product $\alpha h_c$, and say whether a trigger at $h_T=300$ cm can fire.}\sol $m=0.471$; $E=2n+\ell(n-1)=2(1.890)+0.5(0.890)=4.225$ \ans $m^2/\kappa=0.222/(3.5\times10^{-10})=6.34e+08$ \ans $h_c=\frac{1}{0.0750}(6.34e+08)^{1/4.225}=1615$ cm \ans $\alpha h_c=121$ \ans A trigger of 300 cm is below the reach of 1615 cm: it can fire and irrigation will follow \ans

\exbox{Ten-soil example 3c: Loam (S06)}{Loam: $\alpha=0.0360$ cm$^{-1}$, $n=1.560$, $\ell=0.5$, $\kappa=3.5\times10^{-10}$. Compute the exponent $E$, the reach $h_c$ and the product $\alpha h_c$, and say whether a trigger at $h_T=300$ cm can fire.}\sol $m=0.359$; $E=2n+\ell(n-1)=2(1.560)+0.5(0.560)=3.400$ \ans $m^2/\kappa=0.129/(3.5\times10^{-10})=3.68e+08$ \ans $h_c=\frac{1}{0.0360}(3.68e+08)^{1/3.400}=9186$ cm \ans $\alpha h_c=331$ \ans A trigger of 300 cm is below the reach of 9186 cm: it can fire and irrigation will follow \ans

\exbox{Ten-soil example 3c: Silt loam (S07)}{Silt loam: $\alpha=0.0200$ cm$^{-1}$, $n=1.410$, $\ell=0.5$, $\kappa=3.5\times10^{-10}$. Compute the exponent $E$, the reach $h_c$ and the product $\alpha h_c$, and say whether a trigger at $h_T=300$ cm can fire.}\sol $m=0.291$; $E=2n+\ell(n-1)=2(1.410)+0.5(0.410)=3.025$ \ans $m^2/\kappa=0.085/(3.5\times10^{-10})=2.42e+08$ \ans $h_c=\frac{1}{0.0200}(2.42e+08)^{1/3.025}=29528$ cm \ans $\alpha h_c=591$ \ans A trigger of 300 cm is below the reach of 29528 cm: it can fire and irrigation will follow \ans

\exbox{Ten-soil example 3c: Sandy clay loam (S08)}{Sandy clay loam: $\alpha=0.0590$ cm$^{-1}$, $n=1.480$, $\ell=0.5$, $\kappa=3.5\times10^{-10}$. Compute the exponent $E$, the reach $h_c$ and the product $\alpha h_c$, and say whether a trigger at $h_T=300$ cm can fire.}\sol $m=0.324$; $E=2n+\ell(n-1)=2(1.480)+0.5(0.480)=3.200$ \ans $m^2/\kappa=0.105/(3.5\times10^{-10})=3.01e+08$ \ans $h_c=\frac{1}{0.0590}(3.01e+08)^{1/3.200}=7559$ cm \ans $\alpha h_c=446$ \ans A trigger of 300 cm is below the reach of 7559 cm: it can fire and irrigation will follow \ans

\exbox{Ten-soil example 3c: Clay loam (S09)}{Clay loam: $\alpha=0.0190$ cm$^{-1}$, $n=1.310$, $\ell=0.5$, $\kappa=3.5\times10^{-10}$. Compute the exponent $E$, the reach $h_c$ and the product $\alpha h_c$, and say whether a trigger at $h_T=300$ cm can fire.}\sol $m=0.237$; $E=2n+\ell(n-1)=2(1.310)+0.5(0.310)=2.775$ \ans $m^2/\kappa=0.056/(3.5\times10^{-10})=1.6e+08$ \ans $h_c=\frac{1}{0.0190}(1.6e+08)^{1/2.775}=47608$ cm \ans $\alpha h_c=905$ \ans A trigger of 300 cm is below the reach of 47608 cm: it can fire and irrigation will follow \ans

\exbox{Ten-soil example 3c: Silty clay loam (S10)}{Silty clay loam: $\alpha=0.0100$ cm$^{-1}$, $n=1.230$, $\ell=0.5$, $\kappa=3.5\times10^{-10}$. Compute the exponent $E$, the reach $h_c$ and the product $\alpha h_c$, and say whether a trigger at $h_T=300$ cm can fire.}\sol $m=0.187$; $E=2n+\ell(n-1)=2(1.230)+0.5(0.230)=2.575$ \ans $m^2/\kappa=0.035/(3.5\times10^{-10})=9.99e+07$ \ans $h_c=\frac{1}{0.0100}(9.99e+07)^{1/2.575}=127830$ cm \ans $\alpha h_c=1278$ \ans A trigger of 300 cm is below the reach of 127830 cm: it can fire and irrigation will follow \ans

}{}
\IfFileExists{extra/tenby03.tex}{\subsection{Worked solutions on ten crops (hot climate)}
The ten crops of the library differ in their demand, their roots and their sensitivity. The same type of question is answered for each of them at $ET_0=7.5$ mm\,d$^{-1}$; the data are typical values of the crops (Doorenbos--Kassam and FAO 56).

\exbox{Crop example 3c: Maize}{Maize on the loam S06: the trigger fires at the depletion fraction $p=0.55$. Using the demand $r=ET_c/TAW$ and $t_{\rm trig}=p/r$, find the time from full profile to the trigger and the number of irrigations in a season of 125 days.}\sol \givenline{$p=0.55$}\ans $TAW=1.7\times100=170$ mm; \[r=1.20\times7.5/170=0.0529\ \text{d$^{-1}$}\] \ans \[t_{\rm trig}=0.55/0.0529=10.4\ \text{d}\] \ans Irrigations $\approx125/10.4=12$ \finans

\exbox{Crop example 3c: Wheat}{Wheat on the loam S06: the trigger fires at the depletion fraction $p=0.55$. Using the demand $r=ET_c/TAW$ and $t_{\rm trig}=p/r$, find the time from full profile to the trigger and the number of irrigations in a season of 150 days.}\sol \givenline{$p=0.55$}\ans $TAW=1.7\times120=204$ mm; \[r=1.15\times7.5/204=0.0423\ \text{d$^{-1}$}\] \ans \[t_{\rm trig}=0.55/0.0423=13.0\ \text{d}\] \ans Irrigations $\approx150/13.0=12$ \finans

\exbox{Crop example 3c: Cotton}{Cotton on the loam S06: the trigger fires at the depletion fraction $p=0.65$. Using the demand $r=ET_c/TAW$ and $t_{\rm trig}=p/r$, find the time from full profile to the trigger and the number of irrigations in a season of 180 days.}\sol \givenline{$p=0.65$}\ans $TAW=1.7\times110=187$ mm; \[r=1.15\times7.5/187=0.0461\ \text{d$^{-1}$}\] \ans \[t_{\rm trig}=0.65/0.0461=14.1\ \text{d}\] \ans Irrigations $\approx180/14.1=13$ \finans

\exbox{Crop example 3c: Tomato}{Tomato on the loam S06: the trigger fires at the depletion fraction $p=0.40$. Using the demand $r=ET_c/TAW$ and $t_{\rm trig}=p/r$, find the time from full profile to the trigger and the number of irrigations in a season of 135 days.}\sol \givenline{$p=0.40$}\ans $TAW=1.7\times70=119$ mm; \[r=1.15\times7.5/119=0.0725\ \text{d$^{-1}$}\] \ans \[t_{\rm trig}=0.40/0.0725=5.5\ \text{d}\] \ans Irrigations $\approx135/5.5=24$ \finans

\exbox{Crop example 3c: Potato}{Potato on the loam S06: the trigger fires at the depletion fraction $p=0.35$. Using the demand $r=ET_c/TAW$ and $t_{\rm trig}=p/r$, find the time from full profile to the trigger and the number of irrigations in a season of 130 days.}\sol \givenline{$p=0.35$}\ans $TAW=1.7\times50=85$ mm; $r=1.15\times7.5/85=0.1015$ d$^{-1}$ \ans \[t_{\rm trig}=0.35/0.1015=3.4\ \text{d}\] \ans Irrigations $\approx130/3.4=38$ \finans

\exbox{Crop example 3c: Soybean}{Soybean on the loam S06: the trigger fires at the depletion fraction $p=0.50$. Using the demand $r=ET_c/TAW$ and $t_{\rm trig}=p/r$, find the time from full profile to the trigger and the number of irrigations in a season of 130 days.}\sol \givenline{$p=0.50$}\ans $TAW=1.7\times60=102$ mm; \[r=1.15\times7.5/102=0.0846\ \text{d$^{-1}$}\] \ans \[t_{\rm trig}=0.50/0.0846=5.9\ \text{d}\] \ans Irrigations $\approx130/5.9=22$ \finans

\exbox{Crop example 3c: Sugar beet}{Sugar beet on the loam S06: the trigger fires at the depletion fraction $p=0.55$. Using the demand $r=ET_c/TAW$ and $t_{\rm trig}=p/r$, find the time from full profile to the trigger and the number of irrigations in a season of 180 days.}\sol \givenline{$p=0.55$}\ans $TAW=1.7\times90=153$ mm; \[r=1.20\times7.5/153=0.0588\ \text{d$^{-1}$}\] \ans \[t_{\rm trig}=0.55/0.0588=9.4\ \text{d}\] \ans Irrigations $\approx180/9.4=19$ \finans

\exbox{Crop example 3c: Sunflower}{Sunflower on the loam S06: the trigger fires at the depletion fraction $p=0.45$. Using the demand $r=ET_c/TAW$ and $t_{\rm trig}=p/r$, find the time from full profile to the trigger and the number of irrigations in a season of 125 days.}\sol \givenline{$p=0.45$}\ans $TAW=1.7\times80=136$ mm; \[r=1.10\times7.5/136=0.0607\ \text{d$^{-1}$}\] \ans \[t_{\rm trig}=0.45/0.0607=7.4\ \text{d}\] \ans Irrigations $\approx125/7.4=17$ \finans

\exbox{Crop example 3c: Onion}{Onion on the loam S06: the trigger fires at the depletion fraction $p=0.30$. Using the demand $r=ET_c/TAW$ and $t_{\rm trig}=p/r$, find the time from full profile to the trigger and the number of irrigations in a season of 150 days.}\sol \givenline{$p=0.30$}\ans $TAW=1.7\times35=60$ mm; $r=1.05\times7.5/60=0.1324$ d$^{-1}$ \ans \[t_{\rm trig}=0.30/0.1324=2.3\ \text{d}\] \ans Irrigations $\approx150/2.3=66$ \finans

\exbox{Crop example 3c: Alfalfa}{Alfalfa on the loam S06: the trigger fires at the depletion fraction $p=0.55$. Using the demand $r=ET_c/TAW$ and $t_{\rm trig}=p/r$, find the time from full profile to the trigger and the number of irrigations in a season of 100 days.}\sol \givenline{$p=0.55$}\ans $TAW=1.7\times100=170$ mm; \[r=1.15\times7.5/170=0.0507\ \text{d$^{-1}$}\] \ans \[t_{\rm trig}=0.55/0.0507=10.8\ \text{d}\] \ans Irrigations $\approx100/10.8=9$ \finans

}{}
\IfFileExists{extra/tenx03.tex}{\subsection{Worked solutions on the ten soils (second scenario)}
The same questions are repeated for different working conditions, to show how the answer moves with them.

\exbox{Ten-soil example 3b: Coarse loam (S01)}{Coarse loam: $\alpha=0.0961$ cm$^{-1}$, $n=1.353$, $\ell=0.5$, $\kappa=3.5\times10^{-10}$. Compute the exponent $E$, the reach $h_c$ and the product $\alpha h_c$, and say whether a trigger at $h_T=300$ cm can fire.}\sol \givenline{$\alpha=0.0961$ cm$^{-1}$ ; $n=1.353$ ; $\ell=0.5$ ; $\kappa=3.5\times10^{-10}$ ; $h_T=300$ cm can fire}\ans $m=0.261$; $E=2n+\ell(n-1)=2(1.353)+0.5(0.353)=2.882$ \ans $m^2/\kappa=0.068/(3.5\times10^{-10})=1.94e+08$ \ans $h_c=\frac{1}{0.0961}(1.94e+08)^{1/2.882}=7814$ cm \ans $\alpha h_c=751$ \ans A trigger of 300 cm is below the reach of 7814 cm: it can fire and irrigation will follow \finans

\exbox{Ten-soil example 3b: Coarse sand (S02)}{Coarse sand: $\alpha=0.0536$ cm$^{-1}$, $n=1.600$, $\ell=0.5$, $\kappa=3.5\times10^{-10}$. Compute the exponent $E$, the reach $h_c$ and the product $\alpha h_c$, and say whether a trigger at $h_T=300$ cm can fire.}\sol \givenline{$\alpha=0.0536$ cm$^{-1}$ ; $n=1.600$ ; $\ell=0.5$ ; $\kappa=3.5\times10^{-10}$ ; $h_T=300$ cm can fire}\ans $m=0.375$; $E=2n+\ell(n-1)=2(1.600)+0.5(0.600)=3.500$ \ans $m^2/\kappa=0.141/(3.5\times10^{-10})=4.02e+08$ \ans $h_c=\frac{1}{0.0536}(4.02e+08)^{1/3.500}=5359$ cm \ans $\alpha h_c=287$ \ans A trigger of 300 cm is below the reach of 5359 cm: it can fire and irrigation will follow \finans

\exbox{Ten-soil example 3b: Sand (S03)}{Sand: $\alpha=0.1450$ cm$^{-1}$, $n=2.680$, $\ell=0.5$, $\kappa=3.5\times10^{-10}$. Compute the exponent $E$, the reach $h_c$ and the product $\alpha h_c$, and say whether a trigger at $h_T=300$ cm can fire.}\sol \givenline{$\alpha=0.1450$ cm$^{-1}$ ; $n=2.680$ ; $\ell=0.5$ ; $\kappa=3.5\times10^{-10}$ ; $h_T=300$ cm can fire}\ans $m=0.627$; $E=2n+\ell(n-1)=2(2.680)+0.5(1.680)=6.200$ \ans $m^2/\kappa=0.393/(3.5\times10^{-10})=1.12e+09$ \ans $h_c=\frac{1}{0.1450}(1.12e+09)^{1/6.200}=199$ cm \ans $\alpha h_c=29$ \ans A trigger of 300 cm is beyond the reach of 199 cm: the soil cannot supply the demand before it, so the crop stresses first \finans

\exbox{Ten-soil example 3b: Loamy sand (S04)}{Loamy sand: $\alpha=0.1240$ cm$^{-1}$, $n=2.280$, $\ell=0.5$, $\kappa=3.5\times10^{-10}$. Compute the exponent $E$, the reach $h_c$ and the product $\alpha h_c$, and say whether a trigger at $h_T=300$ cm can fire.}\sol \givenline{$\alpha=0.1240$ cm$^{-1}$ ; $n=2.280$ ; $\ell=0.5$ ; $\kappa=3.5\times10^{-10}$ ; $h_T=300$ cm can fire}\ans $m=0.561$; $E=2n+\ell(n-1)=2(2.280)+0.5(1.280)=5.200$ \ans $m^2/\kappa=0.315/(3.5\times10^{-10})=9e+08$ \ans $h_c=\frac{1}{0.1240}(9e+08)^{1/5.200}=425$ cm \ans $\alpha h_c=53$ \ans A trigger of 300 cm is below the reach of 425 cm: it can fire and irrigation will follow \finans

\exbox{Ten-soil example 3b: Sandy loam (S05)}{Sandy loam: $\alpha=0.0750$ cm$^{-1}$, $n=1.890$, $\ell=0.5$, $\kappa=3.5\times10^{-10}$. Compute the exponent $E$, the reach $h_c$ and the product $\alpha h_c$, and say whether a trigger at $h_T=300$ cm can fire.}\sol \givenline{$\alpha=0.0750$ cm$^{-1}$ ; $n=1.890$ ; $\ell=0.5$ ; $\kappa=3.5\times10^{-10}$ ; $h_T=300$ cm can fire}\ans $m=0.471$; $E=2n+\ell(n-1)=2(1.890)+0.5(0.890)=4.225$ \ans $m^2/\kappa=0.222/(3.5\times10^{-10})=6.34e+08$ \ans $h_c=\frac{1}{0.0750}(6.34e+08)^{1/4.225}=1615$ cm \ans $\alpha h_c=121$ \ans A trigger of 300 cm is below the reach of 1615 cm: it can fire and irrigation will follow \finans

\exbox{Ten-soil example 3b: Loam (S06)}{Loam: $\alpha=0.0360$ cm$^{-1}$, $n=1.560$, $\ell=0.5$, $\kappa=3.5\times10^{-10}$. Compute the exponent $E$, the reach $h_c$ and the product $\alpha h_c$, and say whether a trigger at $h_T=300$ cm can fire.}\sol \givenline{$\alpha=0.0360$ cm$^{-1}$ ; $n=1.560$ ; $\ell=0.5$ ; $\kappa=3.5\times10^{-10}$ ; $h_T=300$ cm can fire}\ans $m=0.359$; $E=2n+\ell(n-1)=2(1.560)+0.5(0.560)=3.400$ \ans $m^2/\kappa=0.129/(3.5\times10^{-10})=3.68e+08$ \ans $h_c=\frac{1}{0.0360}(3.68e+08)^{1/3.400}=9186$ cm \ans $\alpha h_c=331$ \ans A trigger of 300 cm is below the reach of 9186 cm: it can fire and irrigation will follow \finans

\exbox{Ten-soil example 3b: Silt loam (S07)}{Silt loam: $\alpha=0.0200$ cm$^{-1}$, $n=1.410$, $\ell=0.5$, $\kappa=3.5\times10^{-10}$. Compute the exponent $E$, the reach $h_c$ and the product $\alpha h_c$, and say whether a trigger at $h_T=300$ cm can fire.}\sol \givenline{$\alpha=0.0200$ cm$^{-1}$ ; $n=1.410$ ; $\ell=0.5$ ; $\kappa=3.5\times10^{-10}$ ; $h_T=300$ cm can fire}\ans $m=0.291$; $E=2n+\ell(n-1)=2(1.410)+0.5(0.410)=3.025$ \ans $m^2/\kappa=0.085/(3.5\times10^{-10})=2.42e+08$ \ans $h_c=\frac{1}{0.0200}(2.42e+08)^{1/3.025}=29528$ cm \ans $\alpha h_c=591$ \ans A trigger of 300 cm is below the reach of 29528 cm: it can fire and irrigation will follow \finans

\exbox{Ten-soil example 3b: Sandy clay loam (S08)}{Sandy clay loam: $\alpha=0.0590$ cm$^{-1}$, $n=1.480$, $\ell=0.5$, $\kappa=3.5\times10^{-10}$. Compute the exponent $E$, the reach $h_c$ and the product $\alpha h_c$, and say whether a trigger at $h_T=300$ cm can fire.}\sol \givenline{$\alpha=0.0590$ cm$^{-1}$ ; $n=1.480$ ; $\ell=0.5$ ; $\kappa=3.5\times10^{-10}$ ; $h_T=300$ cm can fire}\ans $m=0.324$; $E=2n+\ell(n-1)=2(1.480)+0.5(0.480)=3.200$ \ans $m^2/\kappa=0.105/(3.5\times10^{-10})=3.01e+08$ \ans $h_c=\frac{1}{0.0590}(3.01e+08)^{1/3.200}=7559$ cm \ans $\alpha h_c=446$ \ans A trigger of 300 cm is below the reach of 7559 cm: it can fire and irrigation will follow \finans

\exbox{Ten-soil example 3b: Clay loam (S09)}{Clay loam: $\alpha=0.0190$ cm$^{-1}$, $n=1.310$, $\ell=0.5$, $\kappa=3.5\times10^{-10}$. Compute the exponent $E$, the reach $h_c$ and the product $\alpha h_c$, and say whether a trigger at $h_T=300$ cm can fire.}\sol \givenline{$\alpha=0.0190$ cm$^{-1}$ ; $n=1.310$ ; $\ell=0.5$ ; $\kappa=3.5\times10^{-10}$ ; $h_T=300$ cm can fire}\ans $m=0.237$; $E=2n+\ell(n-1)=2(1.310)+0.5(0.310)=2.775$ \ans $m^2/\kappa=0.056/(3.5\times10^{-10})=1.6e+08$ \ans $h_c=\frac{1}{0.0190}(1.6e+08)^{1/2.775}=47608$ cm \ans $\alpha h_c=905$ \ans A trigger of 300 cm is below the reach of 47608 cm: it can fire and irrigation will follow \finans

\exbox{Ten-soil example 3b: Silty clay loam (S10)}{Silty clay loam: $\alpha=0.0100$ cm$^{-1}$, $n=1.230$, $\ell=0.5$, $\kappa=3.5\times10^{-10}$. Compute the exponent $E$, the reach $h_c$ and the product $\alpha h_c$, and say whether a trigger at $h_T=300$ cm can fire.}\sol \givenline{$\alpha=0.0100$ cm$^{-1}$ ; $n=1.230$ ; $\ell=0.5$ ; $\kappa=3.5\times10^{-10}$ ; $h_T=300$ cm can fire}\ans $m=0.187$; $E=2n+\ell(n-1)=2(1.230)+0.5(0.230)=2.575$ \ans $m^2/\kappa=0.035/(3.5\times10^{-10})=9.99e+07$ \ans $h_c=\frac{1}{0.0100}(9.99e+07)^{1/2.575}=127830$ cm \ans $\alpha h_c=1278$ \ans A trigger of 300 cm is below the reach of 127830 cm: it can fire and irrigation will follow \finans

}{}
\IfFileExists{extra/tenbx03.tex}{\subsection{Worked solutions on ten crops (cool climate)}
The ten crops of the library differ in their demand, their roots and their sensitivity. The same type of question is answered for each of them at $ET_0=4.0$ mm\,d$^{-1}$; the data are typical values of the crops (Doorenbos--Kassam and FAO 56).

\exbox{Crop example 3b: Maize}{Maize on the loam S06: the trigger fires at the depletion fraction $p=0.55$. Using the demand $r=ET_c/TAW$ and $t_{\rm trig}=p/r$, find the time from full profile to the trigger and the number of irrigations in a season of 125 days.}\sol $TAW=1.7\times100=170$ mm; $r=1.20\times4.0/170=0.0282$ d$^{-1}$ \ans $t_{\rm trig}=0.55/0.0282=19.5$ d \ans Irrigations $\approx125/19.5=6$ \ans

\exbox{Crop example 3b: Wheat}{Wheat on the loam S06: the trigger fires at the depletion fraction $p=0.55$. Using the demand $r=ET_c/TAW$ and $t_{\rm trig}=p/r$, find the time from full profile to the trigger and the number of irrigations in a season of 150 days.}\sol $TAW=1.7\times120=204$ mm; $r=1.15\times4.0/204=0.0225$ d$^{-1}$ \ans $t_{\rm trig}=0.55/0.0225=24.4$ d \ans Irrigations $\approx150/24.4=6$ \ans

\exbox{Crop example 3b: Cotton}{Cotton on the loam S06: the trigger fires at the depletion fraction $p=0.65$. Using the demand $r=ET_c/TAW$ and $t_{\rm trig}=p/r$, find the time from full profile to the trigger and the number of irrigations in a season of 180 days.}\sol $TAW=1.7\times110=187$ mm; $r=1.15\times4.0/187=0.0246$ d$^{-1}$ \ans $t_{\rm trig}=0.65/0.0246=26.4$ d \ans Irrigations $\approx180/26.4=7$ \ans

\exbox{Crop example 3b: Tomato}{Tomato on the loam S06: the trigger fires at the depletion fraction $p=0.40$. Using the demand $r=ET_c/TAW$ and $t_{\rm trig}=p/r$, find the time from full profile to the trigger and the number of irrigations in a season of 135 days.}\sol $TAW=1.7\times70=119$ mm; $r=1.15\times4.0/119=0.0387$ d$^{-1}$ \ans $t_{\rm trig}=0.40/0.0387=10.3$ d \ans Irrigations $\approx135/10.3=13$ \ans

\exbox{Crop example 3b: Potato}{Potato on the loam S06: the trigger fires at the depletion fraction $p=0.35$. Using the demand $r=ET_c/TAW$ and $t_{\rm trig}=p/r$, find the time from full profile to the trigger and the number of irrigations in a season of 130 days.}\sol $TAW=1.7\times50=85$ mm; $r=1.15\times4.0/85=0.0541$ d$^{-1}$ \ans $t_{\rm trig}=0.35/0.0541=6.5$ d \ans Irrigations $\approx130/6.5=20$ \ans

\exbox{Crop example 3b: Soybean}{Soybean on the loam S06: the trigger fires at the depletion fraction $p=0.50$. Using the demand $r=ET_c/TAW$ and $t_{\rm trig}=p/r$, find the time from full profile to the trigger and the number of irrigations in a season of 130 days.}\sol $TAW=1.7\times60=102$ mm; $r=1.15\times4.0/102=0.0451$ d$^{-1}$ \ans $t_{\rm trig}=0.50/0.0451=11.1$ d \ans Irrigations $\approx130/11.1=12$ \ans

\exbox{Crop example 3b: Sugar beet}{Sugar beet on the loam S06: the trigger fires at the depletion fraction $p=0.55$. Using the demand $r=ET_c/TAW$ and $t_{\rm trig}=p/r$, find the time from full profile to the trigger and the number of irrigations in a season of 180 days.}\sol $TAW=1.7\times90=153$ mm; $r=1.20\times4.0/153=0.0314$ d$^{-1}$ \ans $t_{\rm trig}=0.55/0.0314=17.5$ d \ans Irrigations $\approx180/17.5=10$ \ans

\exbox{Crop example 3b: Sunflower}{Sunflower on the loam S06: the trigger fires at the depletion fraction $p=0.45$. Using the demand $r=ET_c/TAW$ and $t_{\rm trig}=p/r$, find the time from full profile to the trigger and the number of irrigations in a season of 125 days.}\sol $TAW=1.7\times80=136$ mm; $r=1.10\times4.0/136=0.0324$ d$^{-1}$ \ans $t_{\rm trig}=0.45/0.0324=13.9$ d \ans Irrigations $\approx125/13.9=9$ \ans

\exbox{Crop example 3b: Onion}{Onion on the loam S06: the trigger fires at the depletion fraction $p=0.30$. Using the demand $r=ET_c/TAW$ and $t_{\rm trig}=p/r$, find the time from full profile to the trigger and the number of irrigations in a season of 150 days.}\sol $TAW=1.7\times35=60$ mm; $r=1.05\times4.0/60=0.0706$ d$^{-1}$ \ans $t_{\rm trig}=0.30/0.0706=4.2$ d \ans Irrigations $\approx150/4.2=35$ \ans

\exbox{Crop example 3b: Alfalfa}{Alfalfa on the loam S06: the trigger fires at the depletion fraction $p=0.55$. Using the demand $r=ET_c/TAW$ and $t_{\rm trig}=p/r$, find the time from full profile to the trigger and the number of irrigations in a season of 100 days.}\sol $TAW=1.7\times100=170$ mm; $r=1.15\times4.0/170=0.0271$ d$^{-1}$ \ans $t_{\rm trig}=0.55/0.0271=20.3$ d \ans Irrigations $\approx100/20.3=5$ \ans

}{}
\IfFileExists{extra/ex03.tex}{\backhead{EXERCISES}{khred}
\begin{exlist}
\item For the sandy loam S05 ($\alpha=0.075$, $n=1.89$), compute the reach with $\kappa=3.5\times10^{-10}$, $\ell=0.5$.\exans{1615 cm}
\item Is a trigger of 600 cm within the reach of the clay loam S09? Give its reach.\exans{yes; 47608 cm}
\item Find the depletion $f$ at 500 cm for the loam S06.\exans{0.614}
\item Find the suction at 40 per cent depletion in the loam S06 by bisection.\exans{246 cm}
\item A crop in a silt loam (AWC 2.0 mm per cm) roots to 80 cm and is irrigated at depletion 0.45. Find the net dose and the interval at ET = 4.5 mm\,d$^{-1}$.\exans{72 mm; 16.0 d}
\item With the dose law $\phi_D=0.40(1-e^{-d/12})$, find $\phi_D$ for $d=8$, 16 and 24 mm and the water drained in a season of 10 events of 16 mm.\exans{0.195; 0.295; 0.346; 47 mm}
\item The shares of the fate of an irrigation are: transpiration 0.38, soil evaporation 0.24, other evaporation 0.05, drainage 0.28. Find the storage share and the water stored if $I=160$ mm.\exans{0.05; 8 mm}
\item A saving of 40 mm in evaporation is made in year 1 with a memory fate $\phi_D=0.45$ for drainage. How much of the saving appears as drainage?\exans{18 mm}
\item A trigger is set at 900 cm in a soil whose reach is 650 cm. How many irrigations occur? Explain.\exans{none (rainfed)}
\item A tensiometer at 25 cm reads 32 kPa. Convert to suction in cm and pF.\exans{326 cm; pF 2.51}
\end{exlist}
}{}
\IfFileExists{extra/dyk03.tex}{\dykband{DO YOU KNOW ?}
\begin{exlist}
\item Name four classes of irrigation scheduling methods and one weakness of each.
\item What is the difference between a suction trigger and a depletion trigger? Why does the same suction mean different depletions in different soils?
\item Define the reach of a trigger and state what happens when the trigger is beyond it.
\item Why is the constant $\kappa$ the same for all soils in the reach law?
\item What is the marginal fate of an irrigation millimetre?
\item Explain why a small dose per event loses less to drainage.
\item What is hydraulic memory, and why does it limit the banking of a saving?
\item How would you choose a trigger for a coarse soil and for a fine soil?
\end{exlist}
}{}
\IfFileExists{extra/obj03.tex}{\dykband{OBJECTIVE TYPE QUESTIONS}
\begin{exlist}
\item A tensiometer reads\opt{water content}{suction}{salinity}{temperature}
\item A suction trigger of 300 cm in a sand is\opt{a light depletion}{a near-complete depletion}{a saturated soil}{a wilting soil}
\item The reach formula uses the relative conductivity\opt{at saturation}{at the reach equal to $\kappa$}{at wilting only}{of the crop}
\item A trigger beyond the reach\opt{fires often}{never fires}{fires once}{fires at night}
\item The drained share increases with\opt{smaller doses}{larger doses}{the number of events only}{evaporation only}
\item The marginal fates of all fluxes sum to\opt{zero}{one}{the dose}{the reach}
\item Water balance (budget) scheduling needs\opt{a sensor}{the weather and the applied water}{a lysimeter}{satellites}
\item The memory fate is large in soils that\opt{drain freely}{hold water}{are saline}{are frozen}
\item A calendar schedule\opt{responds to the crop}{ignores the state of the crop}{reads a sensor}{needs no water}
\item A granular matrix sensor measures\opt{electrical resistance related to suction}{gamma attenuation}{neutron counts}{leaf temperature}
\end{exlist}
\objkey{1.~(b)\quad 2.~(b)\quad 3.~(b)\quad 4.~(b)\quad 5.~(b)\quad 6.~(b)\quad 7.~(b)\quad 8.~(b)\quad 9.~(b)\quad 10.~(a)}
}{}
